\subsection{Extremal points and barycenters}

PROPOSITION 3. --- Let K be a compact convex subset of a Hausdorff locally convex space E, x a point of K. Every measure \(\mu\) on K, positive, of total mass 1, and admitting x as barycenter, is in the vague closure of the set of discrete positive measures of total mass 1 that admit x as barycenter.

Let U be a neighborhood of \(\mu\) for the vague topology; we can suppose that U consists of the measures \(\nu\) on K such that

\[
\left| \mu (f _ {i}) - \nu (f _ {i}) \right| \leqslant \delta\tag{1}
\]

for a finite number of functions \(f_{i} \in \mathcal{C}(\mathrm{K}; \mathbf{C})\) ( \(1 \leqslant i \leqslant p\) ) and a number \(\delta > 0\) . For every point \(a \in K\) , there exists a closed convex neighborhood \(V_{a}\) of 0 in E such that

\[
\left| f _ {i} (\mathbf {y}) - f _ {i} (\mathbf {a}) \right| \leqslant \delta / 2\tag{2}
\]

for \(1 \leqslant i \leqslant p\) and for every \(\mathbf{y} \in \mathrm{W}_{\mathbf{a}} = \mathrm{K} \cap (\mathbf{a} + \mathrm{V}_{\mathbf{a}})\) . Since \(\mathbf{K}\) is compact, there exists a finite number of points \(\mathbf{a}_j\) ( \(1 \leqslant j \leqslant r\) ) of \(\mathbf{K}\) such that the \(\mathrm{W}_{\mathbf{a}_j}\) form a covering of \(\mathbf{K}\) ( \(1 \leqslant j \leqslant r\) ). Consider a continuous partition of unity \((g_j)_{1 \leqslant j \leqslant r}\) on \(\mathbf{K}\) , subordinate to the covering \((\mathrm{W}_{\mathbf{a}_j})\) , and set \(\alpha_j = \mu(g_j)\) for all \(j\) ; if \(\alpha_j \neq 0\) set \(\mu_j = \alpha_j^{-1} g_j \cdot \mu\) , and if \(\alpha_j = 0\) set \(\mu_j = \varepsilon_{\mathbf{a}_j}\) . Each of the measures \(\mu_j\) is positive, of total mass 1, and its support is contained in the compact convex set \(\mathrm{W}_{\mathbf{a}_j}\) ; moreover, by definition,

\[
\mu = \sum_ {j = 1} ^ {r} \alpha_ {j} \mu_ {j}\tag{3}
\]

since \(g_{j} \cdot \mu = 0\) if \(\mu(g_{j}) = 0\) ; the \(\alpha_{j}\) are \(\geqslant 0\) and

\[
\sum_ {j = 1} ^ {r} \alpha_ {j} = \sum_ {j = 1} ^ {r} \mu (g _ {j}) = \mu \left(\sum_ {j = 1} ^ {r} g _ {j}\right) = \mu (1) = 1.
\]

Let \(\mathbf{x}_j\) be the barycenter of \(\mu_j\) , which belongs to \(\mathrm{W}_{\mathbf{a}_j}\) (No.~1, Prop. 1), and consider the discrete measure \(\nu = \sum_{j=1}^{r} \alpha_j \varepsilon_{\mathbf{x}_j}\) ; it is positive and of total mass 1, and its barycenter is \(\sum_{j=1}^{r} \alpha_j \mathbf{x}_j\) , which is also the barycenter of \(\mu\) by virtue of (3), thus is equal to \(\mathbf{x}\) . Moreover, by (2), \(|f_i(\mathbf{y}) - f_i(\mathbf{a}_j)| \leqslant \delta/2\) for all \(\mathbf{y} \in \mathrm{W}_{\mathbf{a}_j}\) and for all \(i\) , whence, since \(\operatorname{Supp}(\mu_j) \subset \mathrm{W}_{\mathbf{a}_j}\) , \(|\mu_j(f_i) - f_i(\mathbf{a}_j)| \leqslant \delta/2\) for \(1 \leqslant i \leqslant p\) . On the other hand, since \(\mathbf{x}_j \in \mathrm{W}_{\mathbf{a}_j}\) , one also has

\[
\left| \varepsilon_ {\mathbf {x} _ {j}} (f _ {i}) - f _ {i} (\mathbf {a} _ {j}) \right| \leqslant \delta / 2
\]

for \(1 \leqslant i \leqslant p\) , whence \(|\mu_j(f_i) - \varepsilon_{\mathbf{x}_j}(f_i)| \leqslant \delta\) for all \(i\) and \(j\) . Since the \(\alpha_j\) are \(\geqslant 0\) and have sum 1, it follows from (3) and the definition of \(\nu\) that \(\nu\) satisfies the inequality (1).

Q.E.D.

COROLLARY. --- Let \(K'\) be a compact subset of K such that K is the closed convex envelope of \(K'\) . In order that \(x \in K'\) be an extremal point of K, it is necessary and sufficient that \(\varepsilon_{x}\) be the only positive measure on \(K'\) , of total mass 1, having x as barycenter.

Suppose \(\mathbf{x}\) is an extremal point of \(K\) ; to prove that \(\varepsilon_{\mathbf{x}}\) is the only positive measure on \(K'\) , of total mass 1, having \(\mathbf{x}\) as barycenter, it suffices, by Prop. 3, to see that the set of discrete measures \(\nu\) on \(K'\) that are positive, of total mass 1, and have \(\mathbf{x}\) as barycenter, reduces to \(\varepsilon_{\mathbf{x}}\) . But such a measure \(\nu\) is of the form \(\sum_{i=1}^{r} \lambda_i \varepsilon_{\mathbf{x}_i}\) with \(\lambda_i > 0\) for \(1 \leqslant i \leqslant r\) and \(\sum_{i=1}^{r} \lambda_i = 1\) , and the hypothesis that \(\mathbf{x}\) is the barycenter of \(\nu\) may be written \(\mathbf{x} = \sum_{i=1}^{r} \lambda_i \mathbf{x}_i\) . Since \(\mathbf{x}\) is extremal, this implies that \(\mathbf{x}_i = \mathbf{x}\) for all \(i\) , whence \(\nu = \varepsilon_{\mathbf{x}}\) .

Conversely, let us assume that \(\varepsilon_{x}\) is the only positive measure on \(K'\) , of total mass 1, having x as barycenter, and let us show that x is extremal. In the contrary case, there would exist two distinct points \(x', x''\) of K and a real number \(\lambda\) such that \(0 < \lambda < 1\) and \(\mathbf{x} = \lambda\mathbf{x}' + (1 - \lambda)\mathbf{x}''\) . By Prop. 1, \(x'\) (resp. \(x''\) ) is the barycenter of a positive measure \(\mu'\) (resp. \(\mu''\) ) on \(K'\) of total mass 1. Then x is the barycenter of \(\lambda\mu' + (1 - \lambda)\mu''\) . Therefore \(\lambda\mu' + (1 - \lambda)\mu'' = \varepsilon_{x}\) . Therefore \(\mu'\) and \(\mu''\) are proportional to \(\varepsilon_{x}\) , whence \(x' = x'' = x\) , which is absurd.

THEOREM 1 (Choquet). --- Let E be a Hausdorff locally convex space over R, K a metrizable compact convex subset of E, and M the set of extremal points of K. The set M is the intersection of a countable family of open sets in K, and every point of K is the barycenter of a measure \(\mu\) on K such that \(\mu(K - M) = 0\) .

To prove the first assertion, denote by I the interval {[}0,1{]} of \(\mathbf{R}\) ; since K is compact and metrizable, so is \(\mathrm{K} \times \mathrm{K} \times \mathrm{I}\) ; the subset U of \(\mathrm{K} \times \mathrm{K} \times \mathrm{I}\) formed by the triples \((\mathbf{x}, \mathbf{y}, \lambda)\) such that \(\mathbf{x} \neq \mathbf{y}\) and \(0 < \lambda < 1\) is open in \(\mathrm{K} \times \mathrm{K} \times \mathrm{I}\) , therefore there exists a sequence \((\mathrm{F}_n)\) of closed sets in \(\mathrm{K} \times \mathrm{K} \times \mathrm{I}\) whose union is U (GT, IX, §2, No.~5, Prop. 7). The mapping \(q: \mathrm{K} \times \mathrm{K} \times \mathrm{I} \to \mathrm{K}\) defined by \(q(\mathbf{x}, \mathbf{y}, \lambda) = \lambda \mathbf{x} + (1 - \lambda) \mathbf{y}\) is continuous and, by the definition of extremal point, \(\mathrm{K} - \mathrm{M} = q(\mathrm{U}) = \bigcup_{n} q(\mathrm{F}_n)\) ; but \(\mathrm{F}_n\) is compact since it is closed in \(\mathrm{K} \times \mathrm{K} \times \mathrm{I}\) , therefore \(q(\mathrm{F}_n)\) is compact and therefore closed in K; the set \(\mathrm{U}_n = \mathrm{K} - q(\mathrm{F}_n)\) is therefore open in K, and we have \(\mathrm{M} = \bigcap \mathrm{U}_n\) .

In the continuation of the proof, we shall denote by u a continuous and convex numerical function on K, by \(G \subset E \times R\) the graph of u, and by S the closed convex envelope of G in \(E \times R\) .

Lemma 1. --- Let \(\overline{u}\) be the lower envelope of the continuous affine linear functions on E that are \(\geqslant u\) on K. Then S is the set of points \((\mathbf{a}, b) \in \mathrm{E} \times \mathbf{R}\) such that \(\mathbf{a} \in \mathrm{K}\) and \(u(\mathbf{a}) \leqslant b \leqslant \overline{u}(\mathbf{a})\) .

By the Hahn-Banach theorem, for \((\mathbf{a}, b)\) to belong to S, it is necessary and sufficient that \(h(\mathbf{a}, b) \geqslant 0\) for every continuous affine linear function h on \(E \times R\) such that \(h(\mathbf{x}, u(\mathbf{x})) \geqslant 0\) for \(x \in K\) . Setting \(h(\mathbf{x}, t) = f(\mathbf{x}) - \lambda t\) , where f is a continuous affine linear function on E, we see that the relation \((\mathbf{a}, b) \in S\) is equivalent to the following property: the relation

\[
f (\mathbf {x}) \geqslant \lambda u (\mathbf {x}) \quad \text { for   all } \mathbf {x} \in K\tag{4}
\]

implies

\[
f (\mathbf {a}) \geqslant \lambda b.\tag{5}
\]

First set \(\lambda = 0\) ; the fact that (4) implies (5) for every continuous affine linear function \(f\) on \(\mathbf{E}\) is equivalent to the relation \(\mathbf{a} \in \mathbf{K}\) by the Hahn-Banach theorem. Next, set \(\lambda = -1\) and replace \(f\) by \(-f\) ; then the relation \(f(\mathbf{x}) \leqslant u(\mathbf{x})\) in \(\mathbf{K}\) must imply \(f(\mathbf{a}) \leqslant b\) . But since \(u\) is convex and continuous on \(\mathbf{K}\) , \(u(\mathbf{a})\) is equal to the supremum of the \(f(\mathbf{a})\) for the continuous affine linear functions \(f\) on \(\mathbf{E}\) such that \(f(\mathbf{x}) \leqslant u(\mathbf{x})\) on \(\mathbf{K}\) (TVS, II, §5, No.~4, Prop. 5); we thus obtain the relation \(b \geqslant u(\mathbf{a})\) . Finally, set \(\lambda = 1\) ; to say that (4) implies (5) then means, by definition, that \(b \leqslant \overline{u}(\mathbf{a})\) . This proves the lemma, since the relation (4) (resp. (5)) is equivalent to the one obtained by multiplying both sides by a scalar \(>0\) .

Lemma 2. --- If u is strictly convex on K, then \(u(\mathbf{x}) < \overline{u}(\mathbf{x})\) for every non-extremal point x of K.

For, there then exist two distinct points y,z of K such that \(\mathbf{x} = (\mathbf{y} + \mathbf{z})/2\) , whence \(u(\mathbf{x}) < (u(\mathbf{y}) + u(\mathbf{z}))/2\) since u is strictly convex. If f is an affine linear function on E, then \(f(\mathbf{x}) = (f(\mathbf{y}) + f(\mathbf{z}))/2\) ; applying this relation to the continuous affine linear functions f that are \(\geqslant u\) on K, we obtain

\[
\overline {{{u}}} (\mathbf {x}) \geqslant \left(\overline {{{u}}} (\mathbf {y}) + \overline {{{u}}} (\mathbf {z})\right) / 2 \geqslant \left(u (\mathbf {y}) + u (\mathbf {z})\right) / 2 > u (\mathbf {x}).
\]

These lemmas established, we therefore (Lemma 1) have \((\mathbf{a},\overline{u} (\mathbf{a}))\in S\) for all \(\mathbf{a}\in K\) . Since G is compact, being the image of K under the continuous mapping \(\mathbf{x}\mapsto (\mathbf{x},u(\mathbf{x}))\) , there exists, by Prop. 1 of No.~1, a positive measure \(\nu\) on G, of total mass 1, having \((\mathbf{a},\overline{u} (\mathbf{a}))\) as barycenter. Since the restriction of the projection \(\mathrm{pr}_1\) to G is a homeomorphism of G onto K, the measure \(\nu\) may be transported by means of this homeomorphism, which yields a measure \(\mu\) on K (positive and of total mass 1) such that

\[
\mathbf {a} = \int \mathbf {x} d \mu (\mathbf {x}) \quad \mathrm{and} \quad \overline {{{u}}} (\mathbf {a}) = \int u (\mathbf {x}) d \mu (\mathbf {x}).\tag{6}
\]

The first of the relations (6) means that a is the barycenter of \(\mu\) . The function \(\overline{u}\) is upper semi-continuous and bounded on K, hence is \(\mu\) -integrable ( \(\S4\) , No.~4, Cor. 1 of Prop. 5); moreover, since the function \(-\overline{u}\) is by definition convex, we have, by the Cor. of Prop. 2 of No.~1,

\[
\overline {{{u}}} (\mathbf {a}) \geqslant \int \overline {{{u}}} (\mathbf {x}) d \mu (\mathbf {x}),\tag{7}
\]

whence, on comparing with the second of the formulas (6),

\[
\int u (\mathbf {x}) d \mu (\mathbf {x}) \geqslant \int \overline {{{u}}} (\mathbf {x}) d \mu (\mathbf {x}).\tag{8}
\]

But since \(u(\mathbf{x}) \leqslant \overline{u}(\mathbf{x})\) on K, the relation (8) implies that \(u(\mathbf{x}) = \overline{u}(\mathbf{x})\) almost everywhere for \(\mu\) . Taking into account Lemma 2, one sees that Theorem 1 will be proved once the following lemma has been established:

Lemma 3. --- Let E be a Hausdorff locally convex space over R, and K a metrizable compact convex subset of E. Then, there exists a strictly convex numerical function on K.

For, the Banach space \(\mathcal{C}(\mathrm{K};\mathbf{R})\) is separable (GT, X, §3, No.~3, Th. 1), therefore so is the subspace \(\mathcal{A}\) of \(\mathcal{C}(\mathrm{K};\mathbf{R})\) formed by the restrictions to K of the continuous affine linear functions on E. Thus let \((h_n)\) be a sequence dense in \(\mathcal{A}\) , and let \(\alpha_{n} > \sup_{\mathbf{x}\in \mathbb{K}}|h_n(\mathbf{x})|\) . Then each of the functions \(h_n^2 / n^2\alpha_n^2\) is convex in K (TVS, II, §2, No.~8, Examples), and the series with general term \(h_n^2 / n^2\alpha_n^2\) is normally convergent, therefore its sum \(u\) is continuous and convex in K. It remains to see that \(u\) is strictly convex, and for this it suffices to prove that for any two distinct points \(\mathbf{x},\mathbf{x}'\) of K, there is an integer \(n\) such that the restriction of \(h_n^2\) to the segment with endpoints \(\mathbf{x},\mathbf{x}'\) is strictly convex; but for this it suffices that \(h_n(\mathbf{x})\neq h_n(\mathbf{x}')\) (loc. cit.). Now, there exists a function \(h\in \mathcal{A}\) such that \(h(\mathbf{x})\neq h(\mathbf{x}')\) (TVS, II, §4, No.~1, Cor. 1 of Prop. 2) and since the sequence \((h_n)\) is dense in \(\mathcal{A}\) , there exists an \(n\) such that \(h_n(\mathbf{x})\neq h_n(\mathbf{x}')\) .

Q.E.D.

COROLLARY. --- Let E be a Hausdorff locally convex space over R, C a proper convex cone in E with vertex 0, that is complete and metrizable for the uniform structure induced by the weakened uniform structure of E. Let M be the union of the extremal generators of C. For every \(x \in C\) there exist a compact convex subset K of C and a measure \(\lambda \geqslant 0\) on K of total mass 1, such that \(K - (M \cap K)\) is \(\lambda\) -negligible and the barycenter of \(\lambda\) is equal to x.

For, x belongs to a cap K of C (TVS, II, §7, No.~2, Prop. 5), and M∩K contains the set of extremal points of K (loc. cit., Cor. 1 of Prop. 4). It then suffices to apply Th. 1.

\subsection{Applications: I. Vector spaces of continuous real functions}

Let X be a nonempty compact space, H a linear subspace of the Banach space \(\mathcal{C}(\mathrm{X};\mathbf{R})\) that contains the constants and separates the points of X (GT, X, §4, No.~1, Def. 1). We equip \(\mathcal{H}\) with the normed space topology induced by that of \(\mathcal{C}(\mathrm{X};\mathbf{R})\) , and denote by \(\mathcal{H}'\) the dual of this normed space. For every \(x\in \mathrm{X}\) , the mapping \(f\mapsto f(x)\) is a continuous linear form on \(\mathcal{H}\) (the restriction to \(\mathcal{H}\) of the Dirac measure \(\varepsilon_{x}\) ), thus is an element of \(\mathcal{H}'\) that will be denoted \(i_{\mathcal{H}}(x)\) , so that

\[
\langle f, i _ {\mathcal {H}} (x) \rangle = f (x)\tag{9}
\]

for every function \(f \in \mathcal{H}\) and every \(x \in X\) .

The mapping \(i_{\mathcal{H}}\) of X into \(\mathcal{H}'\) is injective and continuous when \(\mathcal{H}'\) is equipped with the weak topology \(\sigma(\mathcal{H}', \mathcal{H})\) ; the second assertion follows at once from the definitions and (9); as for the first, note that if \(x, x'\) are two distinct points of X, by hypothesis there exists a function \(h \in \mathcal{H}\) such that \(h(x) \neq h(x')\) , therefore, by (9), \(\langle h, i_{\mathcal{H}}(x) \rangle \neq \langle h, i_{\mathcal{H}}(x') \rangle\) and a fortiori \(i_{\mathcal{H}}(x) \neq i_{\mathcal{H}}(x')\) . The image \(i_{\mathcal{H}}(\mathrm{X})\) is therefore a compact subset of \(\mathcal{H}'\) (for the weak topology), and \(i_{\mathcal{H}}\) is a homeomorphism of X onto \(i_{\mathcal{H}}(\mathrm{X})\) .

PROPOSITION 4. --- (i) The closed convex envelope C of \(i_{\mathcal{H}}(X)\) in \(\mathcal{H}'\) (for the weak topology \(\sigma(\mathcal{H}', \mathcal{H})\) ) is compact.

\begin{enumerate}
\def\labelenumi{(\roman{enumi})}
\setcounter{enumi}{1}
\tightlist
\item
  For a point \(i_{\mathcal{H}}(x)\) to be an extremal point of \(C\) , it is necessary and sufficient that the only positive measure \(\lambda\) on \(X\) such that
\end{enumerate}

\[
h (x) = \int h d \lambda\tag{10}
\]

for every function \(h \in \mathcal{H}\) (which implies in particular that \(\lambda\) has total mass 1, since \(1 \in \mathcal{H}\) ) be the Dirac measure \(\varepsilon_x\) .

It follows that, for each \(h \in \mathcal{H}\) , the function \(z' \mapsto \langle h, z' \rangle\) on \(C\) attains its supremum at at least one extremal point of \(C\) (TVS, II, §7, No.~1, Prop. 1), and this point belongs to \(i_{\mathcal{H}}(X)\) (loc. cit., Cor. of Prop. 2).

\begin{enumerate}
\def\labelenumi{(\roman{enumi})}
\item
  By (9), \(\| i_{\mathcal{H}}(x)\| \leqslant 1\) in the normed space \(\mathcal{H}'\) , in other words \(i_{\mathcal{H}}(\mathrm{X})\) is bounded, and the assertion follows from the fact that \(\mathcal{H}'\) , equipped with the weak topology \(\sigma (\mathcal{H}',\mathcal{H})\) , is quasi-complete (TVS, III, §4, No.~2, Cor. 5 of Th. 1).
\item
  Every positive measure \(\mu\) of mass 1 on \(i_{\mathcal{H}}(X)\) arises, by transport of structure by means of the homeomorphism \(i_{H}\) , from a positive measure \(\lambda\) of mass 1 on X, the Dirac measure \(\varepsilon_{i_{\mathcal{H}(x)}}\) arising from \(\varepsilon_{x}\) . To say that \(\mu\) admits \(i_{\mathcal{H}(x)}\) as barycenter means, by definition, that
\end{enumerate}

\[
\int_ {\mathrm{X}} \langle h, i _ {\mathcal {H}} (z) \rangle d \lambda (z) = \langle h, i _ {\mathcal {H}} (x) \rangle
\]

for every function \(h \in H\) . Taking (9) into account, the assertion (ii) is just the translation of the criterion of No.~2, Cor. of Prop. 3 for \(i_{\mathcal{H}}(x)\) to be an extremal point of C.

We shall say that a point \(x \in \mathbf{X}\) satisfying condition (ii) of Prop. 4 is \(\mathcal{H}\) -extremal; we denote by \(\mathrm{Ch}_{\mathcal{H}}(\mathbf{X})\) (or simply \(\mathrm{Ch}(\mathbf{X})\) ) the set of these points, and by \(\check{\mathrm{S}}_{\mathcal{H}}(\mathbf{X})\) (or simply \(\check{\mathrm{S}}(\mathbf{X})\) ) the closure of \(\mathrm{Ch}_{\mathcal{H}}(\mathbf{X})\) in \(\mathbf{X}\) .

PROPOSITION 5. --- Every function \(h \in \mathcal{H}\) attains its supremum at at least one \(\mathcal{H}\) -extremal point.

Let \(x\) be a point of \(X\) , \(h\) a function in \(\mathcal{H}\) . The relation \(h(z) \leqslant h(x)\) for all \(z \in X\) may be written \(\langle h, i_{\mathcal{H}}(z) \rangle \leqslant \langle h, i_{\mathcal{H}}(x) \rangle\) for all \(z \in X\) , and therefore means that the weakly closed hyperplane of \(\mathcal{H}'\) with equation \(\langle h, t' \rangle = \langle h, i_{\mathcal{H}}(x) \rangle\) is a support hyperplane of \(i_{\mathcal{H}}(X)\) . It is known (TVS, II, §7, No.~1, Cor. of Prop. 1) that such a hyperplane contains at least one extremal point of the closed convex envelope of \(i_{\mathcal{H}}(X)\) , and such a point \(i_{\mathcal{H}}(y)\) is the image of an \(\mathcal{H}\) -extremal point \(y\) by definition; \(h(y)\) is therefore equal to the supremum of \(h\) in \(X\) .

PROPOSITION 6. --- For every point \(x \in X\) , the following properties are equivalent:

\begin{enumerate}
\def\labelenumi{\alph{enumi})}
\item
  \(x\) is \(\mathcal{H}\) -extremal.
\item
  For every open neighborhood U of x in X and every \(\varepsilon > 0\) , there exists a function \(h \geqslant 0\) in H such that \(h(x) \leqslant \varepsilon\) and \(h(y) \geqslant 1\) for every \(y \in X - U\) .
\end{enumerate}

Let \(x\) be any point of \(X\) , \(f\) a function in \(\mathcal{C}(X; \mathbf{R})\) ; it is known (TVS, II, §3, No.~1, Prop. 1) that the infimum of the numbers \(\lambda(f)\) , for all the positive measures on \(X\) such that \(\lambda(h) = h(x)\) for every function \(h \in \mathcal{H}\) , is equal to the supremum of the numbers \(h(x)\) , where \(h\) runs over the set of functions \(h \in \mathcal{H}\) such that \(h \leqslant f\) . Suppose that \(x\) is \(\mathcal{H}\) -extremal; it then follows from Prop. 4, (ii) that for every function \(f \in \mathcal{C}(X; \mathbf{R})\) ,

\[
f (x) = \sup _ {h \in \mathscr {H}, h \leqslant f} h (x).\tag{11}
\]

To show that \(a)\) implies \(b)\) , we take for \(f\) a continuous mapping of X into \([0,1]\) , with support contained in U, such that \(f(x)=1\) ; then, by (11), there exists a function \(h'\in\mathscr{H}\) such that \(h'\leqslant f\) and \(h'(x)\geqslant 1-\varepsilon\) . Since \(1\in\mathscr{H}\) , the function \(h=1-h'\) meets the conditions of \(b)\) .

Conversely, suppose that the condition \(b)\) is verified; this condition implies that \(1 - h \leqslant \varphi_{\mathrm{U}}\) ; for every positive measure \(\lambda\) on \(X\) satisfying the condition (10), we therefore have

\[
\lambda (\mathrm{U}) = \lambda (\varphi_ {\mathrm{U}}) \geqslant \lambda (1 - h) = 1 - h (x) \geqslant 1 - \varepsilon .
\]

Since, by hypothesis, this relation holds for every \(\varepsilon > 0\) and every open neighborhood U of x, it follows that

\[
\lambda (\{x \}) = \inf _ {U} \lambda (U) \geqslant 1 - \varepsilon
\]

for every \(\varepsilon > 0\) , therefore \(\lambda(\{x\}) = 1\) . Since \(\lambda\) is positive and of total mass 1, necessarily \(\lambda = \varepsilon_x\) , which proves that \(x\) is \(\mathcal{H}\) -extremal, by virtue of Prop. 4, (ii).

PROPOSITION 7. --- Let F be a closed subset of X. The following properties are equivalent:

\begin{enumerate}
\def\labelenumi{\alph{enumi})}
\item
  F contains \(\check{\mathrm{S}}_{\mathcal{H}}(\mathrm{X})\)
\item
  For every function \(h \in \mathcal{H}\) , the set \(F\) intersects the set of points of \(X\) where \(h\) attains its supremum.
\item
  For every point \(x \in X\) , there exists a positive measure \(\mu\) of total mass 1 on \(X\) , such that \(\operatorname{Supp}(\mu) \subset F\) and \(h(x) = \int h d\mu\) for every function \(h \in \mathcal{H}\) .
\end{enumerate}

Let \(G = i_{\mathcal{H}}(F)\) . The condition \(a\) ) signifies that \(G\) contains the set of extremal points of \(C\) . The condition \(b\) ) signifies that \(G\) meets the intersection of \(i_{\mathcal{H}}(X)\) with each of the closed support hyperplanes of \(i_{\mathcal{H}}(X)\) . Finally, the condition \(c\) ) signifies that every point of \(i_{\mathcal{H}}(X)\) is the barycenter of a measure with support contained in \(G\) ; by No.~1, Prop. 1, this is also equivalent to saying that the closed convex envelope of \(i_{\mathcal{H}}(X)\) is equal to the closed convex envelope of \(G\) . The equivalence of the conditions \(a\) ), \(b\) ) and \(c\) ) therefore follows from TVS, II, §7, No.~1, Cor. of Prop. 2.

PROPOSITION 8. --- Suppose X is metrizable. Then the set \(\operatorname{Ch}_{\mathcal{H}}(\mathbf{X})\) of H-extremal points of X is the intersection of a countable family of open sets in X, and for every \(x \in X\) , there exists a positive measure \(\mu\) of total mass 1 on X such that

\[
\mu (\mathrm{X} - \mathrm{Ch} _ {\mathcal {H}} (\mathrm{X})) = 0 \text { and } \int h d \mu = h (x)
\]

for every \(h \in \mathcal{H}\) .

This is the translation of Th. 1 of No.~2, by transport of structure by means of the homeomorphism \(x \mapsto i_{\mathcal{H}}(x)\) , as in Prop. 5.

A certain number of results of this No.~may be extended when H is replaced by a set P of functions defined on X, with values in \(R \cup \{+\infty\}\) , that are lower semi-continuous, P being assumed to contain the constants and to satisfy \(P + P \subset P\) (Exer. 2).

*Example. --- Take X to be the unit ball \(\| \mathbf{x} \| \leqslant 1\) in \(\mathbf{R}^3\) , and let \(\mathcal{H}\) be a vector space of continuous functions on X, containing the restrictions to X of the affine linear functions on \(\mathbf{R}^3\) and satisfying the `maximum principle', that is, for every non-constant function \(h \in \mathcal{H}\) , the set of points of X where \(h\) attains its supremum is contained in the sphere \(\mathbf{S}_2\) . It then follows easily from Props. 5 and 7 that \(\mathrm{Ch}_{\mathcal{H}}(\mathrm{X}) = \check{\mathrm{S}}_{\mathcal{H}}(\mathrm{X}) = \mathbf{S}_2\) . An important example of a vector space \(\mathcal{H}\) satisfying the preceding conditions is the set of functions continuous on X and harmonic in the open ball \(\| \mathbf{x} \| < 1\) . For these functions, one proves that the positive measure \(\mu\) of mass 1 such that \(\operatorname{Supp}(\mu) \subset \mathbf{S}_2\) and \(h(\mathbf{x}) = \int h d\mu\) for all \(h \in \mathcal{H}\) is given, if \(\| \mathbf{x} \| < 1\) , by Poisson's formula

\[
d \mu (\mathbf {z}) = \frac {1 - \| \mathbf {z} \| ^ {2}}{\| \mathbf {z} - \mathbf {x} \| ^ {3}} d \sigma (\mathbf {z}),
\]

where \(\sigma\) is the measure on \(\mathbf{S}_2\) invariant under the orthogonal group and such that \(\sigma(\mathbf{S}_2) = 1\) (Ch. VII, §3, Exer. 8).*

\subsection{Applications: II. Vector spaces of continuous complex functions}

Let X be a nonempty compact space, H a linear subspace of the complex Banach space \(\mathcal{C}(\mathrm{X};\mathbf{C})\) that contains the constants and separates the points of X. The set of real parts \(\mathcal{R}(f)\) of the functions \(f \in H\) is a linear subspace \(H_{r}\) of the real vector space \(\mathcal{C}(\mathrm{X};\mathbf{R})\) ; for every \(f \in H\) , the set \(H_{r}\) also contains \(\mathcal{I}(f) = \mathcal{R}(-if)\) ; it follows that \(H_{r}\) separates the points of X, because the relation \(h(x) = h(y)\) for all \(h \in H_{r}\) implies that \(\mathcal{R}(f(x)) = \mathcal{R}(f(y))\) and \(\mathcal{I}(f(x)) = \mathcal{I}(f(y))\) and so \(f(x) = f(y)\) for all \(f \in H\) . The \(H_{r}\) -extremal points in X are again called H-extremal, the set of them is denoted \(\operatorname{Ch}_{\mathcal{H}}(\mathrm{X})\) , and the closure of the latter set is denoted \(\check{\operatorname{S}}_{\mathcal{H}}(\mathrm{X})\) . The analogues of Props. 5 and 7 are the following:

PROPOSITION 9. --- For every function \(f \in \mathcal{H}\) , \(\mathrm{Ch}_{\mathcal{H}}(\mathrm{X})\) intersects the set of points where \(|f|\) attains its supremum.

We may limit ourselves to the case that \(f\) is not the constant 0. Let \(a\) be a point of \(X\) where \(|f|\) attains its supremum, and set \(g = f / f(a)\) ; then \(g(a) = 1\) and \(|g(x)| \leqslant 1\) for all \(x \in X\) , whence

\[
\mathscr {R} (g (a)) = 1 \quad \text { and } \quad \mathscr {R} (g (x)) \leqslant 1 \quad \text { for   all } x \in X.
\]

By Prop. 5 of No.~3 applied to \(\mathcal{H}_r\) , there exists \(b \in \mathrm{Ch}_{\mathcal{H}}(\mathrm{X})\) where \(\mathcal{R}(g(x))\) attains its supremum 1, whence \(|g(b)| = 1\) since \(|g(b)| \leqslant 1\) ; it follows that \(|f(b)| = |f(a)| \geqslant |f(x)|\) for all \(x \in \mathrm{X}\) .

PROPOSITION 10. --- Let F be a closed subset of X. The following properties are equivalent:

\begin{enumerate}
\def\labelenumi{\alph{enumi})}
\item
  F contains \(\check{\mathrm{S}}_{\mathcal{H}}(\mathrm{X})\)
\item
  For every function \(f \in \mathcal{H}\) , \(F\) intersects the set of points of \(X\) where \(|f|\) attains its supremum.
\item
  For every point \(x \in X\) , there exists a positive measure \(\mu\) of total mass 1 on X such that \(\operatorname{Supp}(\mu) \subset \operatorname{F}\) and \(f(x) = \int f \, d\mu\) for every function \(f \in H\) .
\end{enumerate}

Let us prove the equivalence of the conditions \(a)\) and \(c)\) : let \(f = f_{1} + if_{2}\) with \(f_{1}, f_{2}\) in \(\mathcal{H}_{r}\) ; the relation \(f(x) = \int f d\mu\) is equivalent to the two relations \(f_{1}(x) = \int f_{1} d\mu\) and \(f_{2}(x) = \int f_{2} d\mu\) ; it thus suffices to apply to \(\mathcal{H}_{r}\) the equivalence of the conditions \(a)\) and \(c)\) of Prop. 7 of No.~3. The fact that \(a)\) implies \(b)\) follows from Prop. 9. Let us show that \(b)\) implies \(a)\) ; this is a matter of seeing that if \(b)\) is verified, then, for every \(h \in \mathcal{H}_{r}\) , \(F\) intersects the set of points where \(h\) attains its infimum in \(X\) . The condition \(b)\) implies that \(F\) is nonempty; since \(F\) is compact, there exists \(a \in F\) such that \(h(a) \leqslant h(y)\) for all \(y \in F\) . Let \(f \in \mathcal{H}\) be such that \(h = \mathcal{R}(f)\) ; for every \(\varepsilon > 0\) , the function \(g = f - h(a) + \varepsilon\) belongs to \(\mathcal{H}\) , and

\[
\mathcal {R} (g (y)) = h (y) - h (a) + \varepsilon \geqslant \varepsilon
\]

for all \(y \in F\) . Let c be the supremum of \(|g|\) in X, and set \(b = c^{2}/2\varepsilon\) ; for every \(y \in F\) ,

\[
\left| g (y) - b \right| ^ {2} = \left| g (y) \right| ^ {2} - 2 b \mathcal {R} (g (y)) + b ^ {2} \leqslant c ^ {2} - 2 b \varepsilon + b ^ {2} = b ^ {2},
\]

in other words, the supremum in F of the function \(|g - b|\) is \(\leqslant b\) . Since \(g - b \in \mathcal{H}\) , the hypothesis on F implies that \(|g - b| \leqslant b\) , whence

\[
b ^ {2} \geqslant | g - b | ^ {2} = | g | ^ {2} - 2 b \mathcal {R} (g) + b ^ {2}
\]

and so \(\mathcal{R}(g) \geqslant |g|^{2}/2b \geqslant 0\) ; since \(\mathcal{R}(g) = h - h(a) + \varepsilon\) , and \(\varepsilon > 0\) is arbitrary, we have \(h \geqslant h(a)\) , and \(h(a)\) is the infimum of h in X, which completes the proof.

Remark. --- If f is a continuous real function, a point where \(|f|\) attains its supremum is a point where one of the functions f, -f attains its supremum. For a vector space H of continuous real functions satisfying the hypotheses of No.~3, the Props. 9 and 10 are thus trivial corollaries of Props. 5 and 7, respectively.

\subsection{Applications: III. Algebras of continuous functions}

Lemma 4. --- Let X be a compact space, H a closed linear subspace of the Banach space \(\mathcal{C}(\mathrm{X};\mathbf{C})\) (resp. \(\mathcal{C}(\mathrm{X};\mathbf{R})\) ). Let a be a point of X admitting a countable fundamental system of neighborhoods; assume that, for any numbers c and d such that 0 < c < d < 1 and any open neighborhood U of a, there exists an \(f \in H\) such that

\[
| f | \leqslant 1, \quad | f (a) | \geqslant d, \quad | f (x) | \leqslant c \text {   for   all   } x \in X - U.\tag{12}
\]

Then there exists a function \(u \in \mathcal{H}\) such that \(|u(x)| < |u(a)|\) for all \(x \neq a\) .

Let \(\left(\mathrm{V}_{n}\right)\) ( \(n \geqslant 1\) ) be a fundamental system of neighborhoods of \(a\) , and let \(\lambda, \mu, \varepsilon\) be numbers such that

\[
0 <   \lambda <   1, \quad 1 <   \mu <   \mu + \varepsilon \leqslant 1 + \lambda .
\]

Thus \(0 < \lambda / \mu < 1 / \mu < 1\) . We are going to define, by induction on \(n\) ( \(n \geqslant 1\) ), a decreasing sequence \((\mathrm{U}_n)\) of open neighborhoods of \(a\) such that \(\mathrm{U}_n \subset \mathrm{V}_n\) for all \(n\) , and a sequence \((h_n)\) of functions in \(\mathcal{H}\) satisfying the relations

\[
\left| h _ {n} (x) \right| \leqslant \mu \quad \text { for   all } x \in X\tag{\ensuremath{13_n}}
\]

\[
h _ {n} (a) = 1\tag{\ensuremath{14_n}}
\]

\[
\left| h _ {n} (x) \right| \leqslant \lambda \quad \text { for   all } x \in X - U _ {n}\tag{\ensuremath{15_n}}
\]

\[
\left| \sum_ {j = 1} ^ {n} \lambda^ {j} h _ {j} (y) \right| <   \sum_ {j = 1} ^ {n + 1} \lambda^ {j} \quad \text { for   all } y \in X.\tag{\ensuremath{16_n}}
\]

Assume \(h_m\) and \(U_m\) defined for \(1 \leqslant m < n\) , satisfying the four preceding conditions (with \(n\) replaced by \(m\) ); on the other hand, set \(U_0 = X\) . The function \(\sum_{j=1}^{n-1} \lambda^j h_j\) (equal to 0 if \(n = 1\) ) is continuous and takes the value \(\sum_{j=1}^{n-1} \lambda^j\) at the point \(a\) ; therefore there exists an open neighborhood \(U_n\) of \(a\) , contained in \(U_{n-1} \cap V_n\) , such that

\[
\left| \sum_ {j = 1} ^ {n - 1} \lambda^ {j} h _ {j} (y) \right| <   \sum_ {j = 1} ^ {n - 1} \lambda^ {j} + \varepsilon \lambda^ {n} \quad \text { for   all } y \in \mathrm{U} _ {n}.\tag{17}
\]

By hypothesis there exists a function \(f \in \mathcal{H}\) such that

\[
\begin{array}{l} | f (x) | \leqslant 1 \quad \text { for   all } x \in X, \quad | f (a) | \geqslant 1 / \mu , \\ | f (x) | \leqslant \lambda / \mu \quad \text { for } x \in X - U _ {n}. \end{array}
\]

Set \(h_n = f / f(a)\) ; the relations \((13_n), (14_n)\) and \((15_n)\) then hold; set

\[
g = \sum_ {j = 1} ^ {n} \lambda^ {j} h _ {j} = \sum_ {j = 1} ^ {n - 1} \lambda^ {j} h _ {j} + \lambda^ {n} h _ {n}.
\]

By (17) and \((13_{n})\) , for \(y \in \mathrm{U}_n\) we have

\[
| g (y) | <   \sum_ {j = 1} ^ {n - 1} \lambda^ {j} + \varepsilon \lambda^ {n} + \mu \lambda^ {n} \leqslant \sum_ {j = 1} ^ {n + 1} \lambda^ {j},
\]

since \(\varepsilon + \mu \leqslant 1 + \lambda\) ; for \(x \in \mathrm{X} - \mathrm{U}_n\) , we have \(|h_p(x)| \leqslant \lambda\) for \(1 \leqslant p \leqslant n\) , whence also

\[
| g (x) | \leqslant \sum_ {j = 2} ^ {n + 1} \lambda^ {j} <   \sum_ {j = 1} ^ {n + 1} \lambda^ {j},
\]

which completes the proof of \((16_{n})\) .

This being so, the series \(\sum_{n=1}^{\infty} \lambda^n h_n\) is normally convergent in X since \(\lambda < 1\) and \(|h_n(x)| \leqslant \mu\) for all \(n\) and all \(x \in \mathbf{X}\) ; let \(u\) be its sum, which belongs to \(\mathcal{H}\) since \(\mathcal{H}\) is closed. By the relation \((14_n)\) , we have \(u(a) = \sum_{n=1}^{\infty} \lambda^n\) ; on the other hand if \(x \neq a\) , there exists an integer \(n\) such that \(x \notin \mathrm{U}_{n+1}\) ; therefore \(|h_{n+k}(x)| \leqslant \lambda\) for all \(k \geqslant 1\) by the relation \((15_n)\) ; it follows, using \((16_n)\) , that

\[
\begin{array}{l} | u (x) | \leqslant \left| \sum_ {j = 1} ^ {n} \lambda^ {j} h _ {j} (x) \right| + \left| \sum_ {j = n + 1} ^ {\infty} \lambda^ {j} h _ {j} (x) \right| <   \sum_ {j = 1} ^ {n + 1} \lambda^ {j} + \lambda \sum_ {j = n + 1} ^ {\infty} \lambda^ {j} \\ = \sum_ {j = 1} ^ {\infty} \lambda^ {j} = | u (a) |. \end{array}
\]

THEOREM 2 (E. Bishop). --- Let X be a compact space, A a closed subalgebra of the complex Banach algebra \(\mathcal{C}(\mathrm{X};\mathbf{C})\) . Assume that A contains the constants and separates the points of X. Let \(a\) be a point of X; the following conditions are equivalent:

\begin{enumerate}
\def\labelenumi{\alph{enumi})}
\item
  There exists a function \(f \in \mathcal{A}\) such that \(|f(x)| < |f(a)|\) for all \(x \neq a\) .
\item
  The point \(a\) is \(\mathcal{A}\) -extremal and admits a countable fundamental system of neighborhoods.
\end{enumerate}

\(a) \Rightarrow b)\) : Let \(f \in \mathcal{A}\) be such that \(|f(a)| > |f(x)|\) for \(x \neq a\) ; by Prop. 9 of No.~4, \(a\) is an \(\mathcal{A}\) -extremal point. On the other hand, if \(\mathrm{U}_n\) is the set of \(x \in \mathrm{X}\) such that \(|f(x)| > |f(a)| - 1/n\) , then \(\mathrm{U}_n\) is an open neighborhood of \(a\) , and the intersection of the \(\mathrm{U}_n\) reduces to \(a\) ; since \(\mathrm{X}\) is compact, the \(\mathrm{U}_n\) form a fundamental system of neighborhoods of \(a\) (GT, I, §9, No.~1, Th. 1).

\(b) \Rightarrow a)\) : It suffices to verify that \(b)\) implies the hypotheses of Lemma 4. With the notations of that lemma, set \(\varepsilon = \log d / \log c\) ; thus \(0 < \varepsilon < 1\) . Since \(a\) is an \(\mathcal{A}_r\) -extremal point, there exists a function \(g \in \dot{\mathcal{A}}\) such that

\[
\mathcal {R} (g) \geqslant 0, \quad \mathcal {R} (g (a)) \leqslant \varepsilon , \quad \mathcal {R} (g (x)) \geqslant 1 \text {for} x \in X - U
\]

(No.~3, Prop. 6, b)). Set \(f = c^g\) ; since \(f\) is the sum of the normally convergent series \(\sum_{n=0}^{\infty} (\log c)^n g^n / n!\) , we have \(f \in \mathcal{A}\) and

\[
| f | \leqslant 1, \quad | f (a) | \geqslant c ^ {\varepsilon} = d, \quad | f (x) | \leqslant c \text { for } x \in X - U.\tag{Q.E.D.}
\]

COROLLARY. --- Suppose in addition that X is metrizable. Then the following properties are equivalent:

\begin{enumerate}
\def\labelenumi{\alph{enumi})}
\item
  \(a\) is an \(\mathcal{A}\) -extremal point of X.
\item
  There exists \(u \in \mathcal{A}\) such that \(|u(x)| < |u(a)|\) for all \(x \neq a\) .
\item
  Let \(\mathfrak{M}\) be the set of subsets \(M\) of \(X\) such that for every function \(f \in \mathcal{A}\) , \(|f|\) attains its supremum in \(X\) at at least one point of \(M\) . Then a belongs to all of the sets \(M \in \mathfrak{M}\) .
\item
  Let \(\mathfrak{N}\) be the set of subsets N of X such that, for every function \(f \in \mathcal{A}\) , \(\mathcal{R}(f)\) attains its supremum in X at at least one point of N. Then a belongs to all of the sets \(N \in \mathfrak{N}\) .
\end{enumerate}

In other words,

\[
\operatorname{Ch} _ {\mathcal {A}} (\mathrm{X}) = \bigcap_ {\mathrm{M} \in \mathfrak {M}} \mathrm{M} = \bigcap_ {\mathrm{N} \in \mathfrak {N}} \mathrm{N}.\tag{18}
\]

Since, in a metrizable space, every point admits a countable fundamental system of neighborhoods, the equivalence of \(a\) and \(b\) follows from Th. 2. Let us show that \(b\) implies \(c\) : indeed, \(a\) is the unique point where \(|u|\) attains its supremum; on the other hand, \(c\) implies \(a\) because, for every \(f \in \mathcal{A}\) , \(\mathrm{Ch}_{\mathcal{A}}(\mathrm{X})\) intersects the set of points where \(|f|\) attains its supremum (No.~4, Prop. 9). The same reasoning, using Prop. 5 of No.~3, shows that \(d\) implies \(a\) ). Finally, to see that \(b\) implies \(d\) , we can restrict ourselves to the case that X does not reduce to the single point \(a\) , therefore \(u(a) \neq 0\) ; the function \(v = u / u(a)\) then belongs to \(\mathcal{A}\) , and we have \(v(a) = 1\) and \(|v(x)| < 1\) for \(x \neq a\) , whence \(\mathcal{R}(v(a)) = 1\) and \(\mathcal{R}(v(x)) < 1\) for \(x \neq a\) . Since the function \(\mathcal{R}(v)\) attains its supremum only at the point \(a\) , we have indeed \(a \in N\) for every \(N \in \mathfrak{N}\) .

*Examples. --- Let \(\mathrm{X}_1\) be the set of points \((z_1, z_2) \in \mathbf{C}^2\) such that \(|z_1|^2 + |z_2|^2 \leqslant 1\) (the unit ball in \(\mathbf{R}^4\) ) and let \(\mathcal{A}_1'\) be the set of restrictions to \(\mathrm{X}_1\) of the holomorphic functions, with values in \(\mathbf{C}\) , defined in a neighborhood of \(\mathrm{X}_1\) in \(\mathbf{C}^2\) (the neighborhood depending on the function considered); let \(\mathcal{A}_1\) be the closure of \(\mathcal{A}_1'\) in \(\mathcal{C}(\mathrm{X}_1; \mathbf{C})\) , which is obviously a closed complex subalgebra of \(\mathcal{C}(\mathrm{X}_1; \mathbf{C})\) and separates the points of \(\mathrm{X}_1\) . Application of the `maximum principle' for holomorphic functions shows that \(\mathrm{Ch}_{\mathcal{A}_1}(\mathrm{X}_1)\) is the sphere \(\mathbf{S}_3\) .

In the preceding definition, let us replace \(X_1\) by the `polydisk' \(X_2\) defined by the relations \(|z_1| \leqslant 1\) and \(|z_2| \leqslant 1\) , which yields subalgebras \(\mathscr{A}_2'\) and \(\mathscr{A}_2\) (the closure of \(\mathcal{A}_2^{\prime}\) ) of \(\mathcal{C}(\mathrm{X}_2; \mathbf{C})\) . Here, the maximum principle shows that \(\mathrm{Ch}_{\mathcal{A}_2}(\mathrm{X}_2)\) is the `torus' defined by the relations \(|z_1| = 1\) and \(|z_2| = 1\) .

From these results, one deduces that there does not exist an analytic isomorphism of an open neighborhood of \(X_{1}\) onto an open neighborhood of \(X_{2}\) that transforms \(X_{1}\) into \(X_{2}\) ; for, if v were the restriction to \(X_{1}\) of such a mapping, one would have \(A_{2}=vA_{1}v^{-1}\) and so v would transform \(S_{3}\) into a space homeomorphic to \(T^{2}\) , which is absurd since \(S_{3}\) is simply connected but \(T^{2}\) is not. One will observe, however, that the spaces \(X_{1}\) and \(X_{2}\) are homeomorphic, both being bounded convex sets in \(R^{4}\) with nonempty interior.

\subsection{Uniqueness of integral representations}

Let E be a Hausdorff weak locally convex space (TVS, II, §6, No.~2), C a proper pointed convex cone in E. One knows that C is the set of elements \(\geqslant0\) of E for an order relation compatible with the vector space structure of E. When C is said to be lattice-ordered, it is of course the order induced on C by that of E that is understood.

Lemma 5. --- Assume that C is weakly complete. Let A be the set of restrictions to C of the continuous linear forms on E. Let \((f_{\lambda})_{\lambda \in \Lambda}\) be a finite family of elements of A, and \(f = \sup(f_{\lambda})\) . For every \(x \in C\) , define

\[
\overline {{{f}}} (x) = \sup \left(f (x _ {1}) + f (x _ {2}) + \dots + f (x _ {n})\right),
\]

the supremum being taken over the set \(S_{x}\) of sequences \((x_{1}, x_{2}, \ldots, x_{n})\) of elements of C such that \(x_{1} + x_{2} + \cdots + x_{n} = x\) . Let \(\operatorname{Card}\Lambda = p\) . Then there exists \((y_{1}, \ldots, y_{p}) \in \operatorname{S}_{x}\) such that \(\overline{f}(x) = f(y_{1}) + \cdots + f(y_{p})\) .

Denote by \(f_1, f_2, \ldots, f_p\) the elements of the family \((f_\lambda)\) . For \(k = 1, 2, \ldots, p\) , let \(C_k\) be the set of \(y \in C\) such that

\[
f _ {1} (y) <   f (y), f _ {2} (y) <   f (y), \dots , f _ {k - 1} (y) <   f (y), f _ {k} (y) = f (y).
\]

The \(C_{k}\) are disjoint convex cones with union C. Let \(x_{1}, x_{2}, \ldots, x_{n}\) in C be such that \(x_{1} + x_{2} + \cdots + x_{n} = x\) . Let \(y_{k}\) be the sum of the \(x_{i}\) that belong to \(C_{k}\) . Then \(y_{1} + y_{2} + \cdots + y_{p} = x\) . Since f is affine on \(C_{k}\) , \(f(y_{1}) + \cdots + f(y_{p}) = f(x_{1}) + \cdots + f(x_{n})\) . Therefore

\[
f (x) = \sup \bigl (f (y _ {1}) + \dots + f (y _ {p}) \bigr),\tag{19}
\]

where \((y_{1}, y_{2}, \ldots, y_{p})\) runs over the set of sequences of p points of C such that \(y_{1} + y_{2} + \cdots + y_{p} = x\) . Set \(D = C \cap (x - C)\) . Since D is compact (TVS, II, §6, No.~8, Cor. 2 of Prop. 11), so is the set of elements \((y_{1}, \ldots, y_{p})\) of \(D^{p}\) such that \(y_{1} + \cdots + y_{p} = x\) , thus the supremum (19) is attained.

Lemma 6. --- We maintain the hypotheses and notations of Lemma 5, and assume the \(f_{\lambda}\) to be positive. The function \(\overline{f}\) is positively homogeneous, concave and upper semi-continuous in C. It is affine if C is lattice-ordered.

It is clear that \(\overline{f}\) is positively homogeneous. Let x, y belong to C. If \(x_{1}, \ldots, x_{m}, y_{1}, \ldots, y_{n}\) in C are such that \(x_{1} + \cdots + x_{m} = x\) , \(y_{1} + \cdots + y_{n} = y\) , then \(x_{1} + \cdots + x_{m} + y_{1} + \cdots + y_{n} = x + y\) , therefore

\[
f (x _ {1}) + \dots + f (x _ {m}) + f (y _ {1}) + \dots + f (y _ {n}) \leqslant \overline {{f}} (x + y);
\]

it follows that \(\overline{f}(x) + \overline{f}(y) \leqslant \overline{f}(x + y)\) , therefore \(\overline{f}\) is concave. Let \(L\) (resp. \(L_{\lambda}\) ) be the set of \((t, x) \in \mathbf{R} \times E\) such that \(x \in C\) and \(0 \leqslant t \leqslant \overline{f}(x)\) (resp. \(0 \leqslant t \leqslant f_{\lambda}(x)\) ). Each of the \(L_{\lambda}\) is closed in the weakly complete proper convex cone \(\mathbf{R}_{+} \times C\) , therefore the sum \(\sum_{\lambda \in \Lambda} L_{\lambda}\) is closed (TVS, II, §6, No.~8, Cor. 2 of Prop. 11). By Lemma 5, this sum is equal to L. Therefore L is closed, which proves that \(\overline{f}\) is upper semicontinuous. Finally, assume C is lattice-ordered, and let us prove that \(\overline{f}\) is convex. Let \(x, y\) belong to C and let \(\varepsilon > 0\) . There exist \(z_1, z_2, \ldots, z_n\) in C such that \(f(z_1) + \cdots + f(z_n) \geqslant \overline{f}(x + y) - \varepsilon\) and \(z_1 + \cdots + z_n = x + y\) . The vector space C - C is lattice-ordered for the order induced by that of E (A, VI, §1, No.~9, Prop. 8). By the decomposition theorem (loc. cit., No.~10, Th. 1), there exist \(x_1, \ldots, x_n, y_1, \ldots, y_n\) in C such that

\[
x _ {1} + y _ {1} = z _ {1}, \dots , x _ {n} + y _ {n} = z _ {n}, x _ {1} + \dots + x _ {n} = x, y _ {1} + \dots + y _ {n} = y.
\]

Then, since \(f\) is positively homogeneous and convex,

\[
\begin{array}{r l} \overline {{f}} (x + y) & \leqslant \varepsilon + f (z _ {1}) + \dots + f (z _ {n}) \\ & \leqslant \varepsilon + f (x _ {1}) + f (y _ {1}) + \dots + f (x _ {n}) + f (y _ {n}) \\ & \leqslant \varepsilon + \overline {{f}} (x) + \overline {{f}} (y). \end{array}
\]

Since \(\varepsilon\) is an arbitrary number \(>0\) , we have proved that \(\overline{f}\) is indeed convex.

THEOREM 3 (Choquet). --- Let E be a Hausdorff weak locally convex space, C a weakly complete proper convex cone with vertex 0 in E, G the union of the extremal generators of C, K a compact convex subset of C, \(\lambda\) and \(\lambda'\) positive measures of mass 1 on K, admitting the same barycenter, such that \(\lambda^{*}(\mathrm{K} - (\mathrm{K} \cap \mathrm{G})) = \lambda'^{*}(\mathrm{K} - (\mathrm{K} \cap \mathrm{G})) = 0\) . Assume that C is lattice-ordered. Then, for every lower semi-continuous, positively homogeneous convex function \(f \geqslant 0\) on C, \(\lambda^{*}(f|K) = \lambda'^{*}(f|K)\) .

Let \(\mathcal{A}\) (resp. \(\mathcal{A}'\) ) be the set of restrictions to \(\dot{\mathrm{C}}\) of the continuous linear forms (resp. affine functions) on E. We know (TVS, II, §5, No.~4, Remark 2) that \(f\) is the upper envelope of the set of elements of \(\mathcal{A}\) that are \(\leqslant f\) . The set of functions of the form \(\sup(f_1, \ldots, f_p)\) , where \(f_1, \ldots, f_p\) belong to \(\mathscr{A}\) , \(f_1 \geqslant 0, \ldots, f_p \geqslant 0\) , is an increasing directed set and has \(f\) as its upper envelope. Taking into account §1, No.~1, Th. 1, it suffices to verify the equality \(\lambda(f|K) = \lambda'(f|K)\) when \(f\) is of the preceding form.

Define \(\overline{f}\) as in Lemma 5. It is clear that \(\overline{f}(y) = f(y)\) if \(y \in G\) . Since \(\lambda^{*}(\mathrm{K} - (\mathrm{K} \cap \mathrm{G})) = 0\) , we have \(\lambda(f|K) = \lambda(\overline{f}|K)\) . By Lemma 6, \(\overline{f}\) is affine and upper semi-continuous. Therefore \(\overline{f}|K\) is the lower envelope of a decreasing directed set of restrictions of elements of \(\mathcal{A}'\) to K (TVS, II, §5, No.~4, Prop. 6). Let \(x \in K\) be the barycenter of \(\lambda\) . If \(g \in \mathcal{A}\) then \(\lambda(g|K) = g(x)\) . Therefore \(\lambda(\overline{f}|K) = \overline{f}(x)\) (§4, No.~4, Cor. 2 of Prop. 5). Thus \(\lambda(f|K) = \overline{f}(x)\) , and one sees similarly that \(\lambda'(f|K) = \overline{f}(x)\) .

COROLLARY. --- Let E be a Hausdorff locally convex space, C a proper convex cone with vertex 0 in E, admitting a compact sole M, and let G be the union of the extremal generators of C. Let \(x \in M\) . If C is lattice-ordered, then there exists at most one positive measure \(\lambda\) of mass 1 on M, such that \(\lambda^{*}(\mathrm{M} - (\mathrm{G} \cap \mathrm{M})) = 0\) , and admitting x as barycenter.

Replacing the topology of E by the weakened topology (which does not change the topology of M), one can suppose E to be a weak space. Let \(\lambda\) and \(\lambda'\) be two measures on M having the stated properties, and let h be a continuous linear form on E such that M is the intersection of C and the hyperplane with equation \(h(x)=1\) . Let S be the subset of \(\mathcal{C}(\mathrm{M})\) consisting of the restrictions to M of the positively homogeneous and continuous convex functions \(\geqslant0\) on C. The cone C is weakly complete (TVS, II, §7, No.~3). By Th. 3, \(\lambda(f)=\lambda'(f)\) for every \(f\in S\) .

If \(f_{1}, f_{2}, f_{3}, f_{4}\) belong to \(\mathcal{S}\) , then

\[
\sup \left(f _ {1} - f _ {2}, f _ {3} - f _ {4}\right) = \sup \left(f _ {1} + f _ {4}, f _ {3} + f _ {2}\right) - \left(f _ {2} + f _ {4}\right) \in \mathscr {S} - \mathscr {S}
\]

\[
\inf \left(f _ {1} - f _ {2}, f _ {3} - f _ {4}\right) = - \sup \left(f _ {2} - f _ {1}, f _ {4} - f _ {3}\right) \in \mathscr {S} - \mathscr {S}.
\]

Since \(h|M \in \mathcal{S}\) , \(\mathcal{S} - \mathcal{S}\) contains the constant functions. If \(x\) and \(y\) are two distinct points of \(M\) , there exists a continuous linear form on \(E\) that separates \(x\) and \(y\) , and this form is the difference of two continuous linear forms that are positive on \(C\) (TVS, II, §6, No.~8, Lemma 1). It follows from the foregoing that for \(\alpha, \beta\) real, there exists \(f \in \mathcal{S} - \mathcal{S}\) such that \(f(x) = \alpha\) , \(f(y) = \beta\) . Then \(\mathcal{S} - \mathcal{S}\) is dense in \(\mathcal{C}(M)\) for the topology of uniform convergence (GT, X, §4, No.~1, Cor. of Prop. 2). Since \(\lambda\) and \(\lambda'\) coincide on \(\mathcal{S} - \mathcal{S}\) , we have \(\lambda = \lambda'\) .

\chapter*{Exercises}
\phantomsection
\addcontentsline{toc}{chapter}{Exercises}
\setcounter{exercisegroup}{1}
\edef\CurrentExerciseGroup{\arabic{exercisegroup}}
\section*{\texorpdfstring{\S\CurrentExerciseGroup}{Section \CurrentExerciseGroup}}
\phantomsection
\addcontentsline{toc}{section}{\texorpdfstring{Exercises for \S\CurrentExerciseGroup}{Exercises for Section \CurrentExerciseGroup}}

\begin{enumerate}
\def\labelenumi{\arabic{enumi})}
\tightlist
\item
  Show that if \(f\) and \(g\) are two numerical functions \(\geqslant 0\) defined on \(X\) , and \(\mu\) is a positive measure on \(X\) , then
\end{enumerate}

\[
\mu^ {*} (\sup (f, g)) + \mu^ {*} (\inf (f, g)) \leqslant \mu^ {*} (f) + \mu^ {*} (g).
\]

From this, deduce that if A and B are any subsets of X, then

\[
\mu^ {*} (\mathrm{A} \cup \mathrm{B}) + \mu^ {*} (\mathrm{A} \cap \mathrm{B}) \leqslant \mu^ {*} (\mathrm{A}) + \mu^ {*} (\mathrm{B})
\]

(cf.~§4, Exer. 8 d)).

\begin{enumerate}
\def\labelenumi{\arabic{enumi})}
\setcounter{enumi}{1}
\item
  For every \(x \in \mathbf{X}\) , show that for every numerical function \(f \geqslant 0\) defined on \(\mathbf{X}\) , \(\varepsilon_x^*(f) = f(x)\) .
\item
  Give an example, in \(\mathbf{R}\) , of an open set that is not relatively compact, whose outer measure (for Lebesgue measure) is finite.
\item
  Let \(\mathrm{X}\) be a locally compact space, \(\alpha\) a finite numerical function \(\geqslant 0\) on \(\mathrm{X}\) such that, for every compact subset \(\mathrm{K}\) of \(\mathrm{X}\) , \(\sum_{x\in \mathrm{K}}\alpha (x)\) is finite; let \(\mu\) be the positive measure
\end{enumerate}

on X defined by the masses \(\alpha(x)\) (Ch. III, §1, No.~3, Example I).

\begin{enumerate}
\def\labelenumi{\alph{enumi})}
\tightlist
\item
  Show that for every function \(f \geqslant 0\) , lower semi-continuous on \(X\) , one has \(\mu^{*}(f) = \sum_{x \in X} \alpha(x)f(x)\) , with the convention that \(\alpha(x)f(x) = 0\) when \(\alpha(x) = 0\) and \(f(x) = -\infty\) .
\end{enumerate}

\[
f (x) = + \infty
\]

\begin{enumerate}
\def\labelenumi{\alph{enumi})}
\setcounter{enumi}{1}
\tightlist
\item
  Let \(f \geqslant 0\) be any numerical function defined on \(X\) . Show that if \(\mu^{*}(f) < +\infty\) , then \(\mu^{*}(f) = \sum_{x \in X} \alpha(x)f(x)\) , with the same convention as in a). (Cf. Exer. 5.)
\end{enumerate}

¶ 5) Let X be the subset of the plane \(\mathbf{R}^2\) formed by the union of the line \(D = \{0\} \times \mathbf{R}\) and the set of points \((1/n, k/n^2)\) , where \(n\) runs over the set of integers \(>0\) , and \(k\) over the set \(\mathbf{Z}\) of rational integers.

\begin{enumerate}
\def\labelenumi{\alph{enumi})}
\item
  For every point \((0,y)\) of D and every integer n \textgreater{} 0, let \(\mathrm{T}_{n}(y)\) be the set of points \((u,v)\) in X such that \(u \leqslant 1/n\) and \(|v - y| \leqslant u\) . Show that if one takes as fundamental system of neighborhoods of each point \((0,y)\) in D the set of \(\mathrm{T}_{n}(y)\) (n \textgreater{} 0), and as fundamental system of neighborhoods of each of the other points of X the unique set reduced to that point, one defines on X a topology T for which X is a locally compact space that is not paracompact.
\item
  Let \(\alpha\) be the numerical function on \(X\) that is equal to 0 on \(D\) and to \(1/n^3\) at each of the points \((1/n, k/n^2)\) . Show that, for every compact subset \(K\) of \(X\) , \(\sum_{x \in K} \alpha(x)\)
\end{enumerate}

is finite; let \(\mu\) be the positive measure on X defined by the masses \(\alpha(x)\) . Show that \(\mu^{*}(D) = +\infty\) even though \(\alpha(x) = 0\) on D. (If an open set U for \(\mathcal{T}\) contains D, show that there exist an interval \(a \leqslant y \leqslant b\) on D not reduced to a point, a set B dense in this interval (for the usual topology of \(\mathbf{R}\) ), and an integer \(n > 0\) such that, for every \(y \in B\) , one has \(\mathrm{T}_{n}(y) \subset \mathrm{U}\) ; for this, one may make use of Baire's theorem.)

\begin{enumerate}
\def\labelenumi{\arabic{enumi})}
\setcounter{enumi}{5}
\tightlist
\item
  Let \(f\) be a numerical function \(\geqslant 0\) defined on \(X\) .
\end{enumerate}

\begin{enumerate}
\def\labelenumi{\alph{enumi})}
\item
  Show that, for the mapping \(\mu \mapsto \mu^{*}(f)\) of \(\mathcal{M}_{+}(X)\) into \(\overline{\mathbf{R}}\) to be continuous for the vague topology, it is necessary (and sufficient) that \(f\) be continuous with compact support (make use of Exer. 2). For \(\mu \mapsto \mu^{*}(f)\) to be lower semi-continuous for the vague topology, it is necessary (and sufficient, cf.~Prop. 4) that \(f\) be lower semi-continuous.
\item
  Show that, for the mapping \(\mu \mapsto \mu^{*}(f)\) of \(\mathcal{M}_{+}(\mathrm{X})\) into \(\overline{\mathbf{R}}\) to be continuous for the quasi-strong topology (Ch. III, §1, Exer. 8), it is necessary and sufficient that \(f\) be bounded and have compact support (method analogous to that of \(a\) )). From this, deduce that for every function \(f \geqslant 0\) that is zero on the complement of a countable union of compact sets, the mapping \(\mu \mapsto \mu^{*}(f)\) is lower semi-continuous for the quasi-strong topology (make use of Th. 3) (cf.~Exer. 7 b)).
\item
  Show that, for the mapping \(\mu \mapsto \mu^{*}(f)\) of \(\mathcal{M}_{+}(\mathrm{X}) \cap \mathcal{M}^{1}(\mathrm{X})\) into \(\overline{\mathbf{R}}\) to be continuous for the ultrastrong topology (Ch. III, §1, Exer. 15), it is necessary and sufficient that \(f\) be bounded; for every function \(f \geqslant 0\) defined on \(\mathrm{X}\) , \(\mu \mapsto \mu^{*}(f)\) is lower semicontinuous for the ultrastrong topology.
\item
  Show that, for the mapping \(\mu \mapsto \mu^{*}(f)\) of \(\mathcal{M}_{+}(\mathrm{X}) \cap \mathcal{M}^{1}(\mathrm{X})\) into \(\overline{\mathbf{R}}\) to be continuous for the weak topology (Ch. III, §1, Exer. 15), it is necessary and sufficient that \(f\) be continuous on \(\mathrm{X}\) and tend to 0 at the point at infinity; for \(\mu \mapsto \mu^{*}(f)\) to be lower semi-continuous for the weak topology, it is necessary and sufficient that \(f\) be lower semi-continuous.
\end{enumerate}

\begin{enumerate}
\def\labelenumi{\arabic{enumi})}
\setcounter{enumi}{6}
\tightlist
\item
  \begin{enumerate}
  \def\labelenumii{\alph{enumii})}
  \tightlist
  \item
    Let \((\mu_{n})\) be an increasing sequence of measures \(\geqslant0\) on a locally compact space X; assume that the sequence is bounded above in \(\mathcal{M}_{+}(X)\) and denote by \(\mu\) its supremum. Let f be a function \(\geqslant0\) defined on X and zero on the complement of a countable union of compact sets. Show that \(\mu^{*}(f)=\sup\mu_{n}^{*}(f)\) (cf.~Exer. 6 b)).
  \end{enumerate}
\end{enumerate}

\begin{enumerate}
\def\labelenumi{\alph{enumi})}
\setcounter{enumi}{1}
\tightlist
\item
  Let X and \(\mu\) be the locally compact space and the measure defined in Exer. 5. Let \(\alpha_{n}\) be the numerical function equal to \(\alpha\) (in the notation of Exer. 5) for every point \((1/m, k/m^{2})\) such that \(m \leqslant n\) , and equal to 0 at the other points of X; let \(\mu_{n}\) be the measure defined by the masses \(\alpha_{n}(x)\) . Show that \(\mu\) is the supremum of the sequence \((\mu_{n})\) in \(\mathcal{M}_{+}(X)\) and that \(\mu_{n}^{*}(D)=0\) for all n, but \(\mu^{*}(D)=+\infty\) .
\end{enumerate}

\begin{enumerate}
\def\labelenumi{\arabic{enumi})}
\setcounter{enumi}{7}
\tightlist
\item
  \begin{enumerate}
  \def\labelenumii{\alph{enumii})}
  \tightlist
  \item
    Let \((\mu_n)\) be a decreasing sequence of measures \(\geqslant 0\) on a locally compact space \(X\) , and let \(\mu\) be the infimum of the sequence in \(\mathcal{M}_+(X)\) . Show that if \(f\) is a function \(\geqslant 0\) such that \(\mu_n^*(f) < +\infty\) from some index onward, then \(\mu^*(f) = \inf_n \mu_n^*(f)\) (when \(g\) is lower semi-continuous, positive, and satisfies \(\mu^*(g) < +\infty\) , note that there exists a sequence \((h_m)\) of continuous functions \(\geqslant 0\) with compact support, such that \(\sum_{m=1}^{\infty} h_m \leqslant g\) and \(\mu_n^*(g) = \sum_{m=1}^{\infty} \mu_n^*(h_m)\) for every index \(n\) .
  \end{enumerate}
\end{enumerate}

\begin{enumerate}
\def\labelenumi{\alph{enumi})}
\setcounter{enumi}{1}
\tightlist
\item
  On the discrete space X = N, let \(\mu_{n}\) be the measure defined by placing mass +1 at each point \(m \geqslant n\) . Show that the infimum of the decreasing sequence \((\mu_{n})\) in \(\mathcal{M}_{+}(X)\) is 0, but that \(\mu_{n}^{*}(X) = +\infty\) for all n.
\end{enumerate}

\setcounter{exercisegroup}{3}
\edef\CurrentExerciseGroup{\arabic{exercisegroup}}
\section*{\texorpdfstring{\S\CurrentExerciseGroup}{Section \CurrentExerciseGroup}}
\phantomsection
\addcontentsline{toc}{section}{\texorpdfstring{Exercises for \S\CurrentExerciseGroup}{Exercises for Section \CurrentExerciseGroup}}

\begin{enumerate}
\def\labelenumi{\arabic{enumi})}
\item
  Let \(\mu\) be Lebesgue measure on the interval \(X = [0,1[\) of \(\mathbf{R}\) . For every integer \(n = 2^h + k\) ( \(h \geqslant 0\) , \(0 \leqslant k < 2^h\) ) let \(f_n\) be the function equal to 1 on the interval \([k/2^h, (k+1)/2^h[\) and to 0 elsewhere in \(X\) . Show that the sequence \((f_n)\) converges in mean of order \(p\) to 0 for every \(p > 0\) , but that the sequence \((f_n(x))\) does not converge for any point \(x \in X\) .
\item
  Show that every numerical function \(f\) belonging to \(\mathcal{L}^p(\mathrm{X};\mathbf{R})\) is equal almost everywhere to the difference \(g_1 - g_2\) of two lower semi-continuous positive functions belonging to \(\mathcal{L}^p(\mathrm{X};\mathbf{R})\) (note that \(f(x)\) is equal almost everywhere to the sum of an absolutely convergent series \(\sum_{n=1}^{\infty} f_n(x)\) , where the \(f_n\) are continuous with compact support).
\item
  Let \(\mathbf{X}\) be a locally compact space, \(\alpha\) a finite numerical function \(\geqslant 0\) defined on \(\mathbf{X}\) such that, for every compact subset \(\mathbf{K}\) of \(\mathbf{X}\) , \(\sum_{x\in \mathbf{K}}\alpha (x) < +\infty\) . Let \(\mu\) be the measure on \(\mathbf{X}\) defined by the masses \(\alpha (x)\) . Show that, for this measure, \(\mathcal{F}_{\mathrm{F}}^{p} = \mathcal{L}_{\mathrm{F}}^{p}\) for every finite \(p\geqslant 1\) and every Banach space \(\mathbf{F}\) .
\end{enumerate}

\setcounter{exercisegroup}{4}
\edef\CurrentExerciseGroup{\arabic{exercisegroup}}
\section*{\texorpdfstring{\S\CurrentExerciseGroup}{Section \CurrentExerciseGroup}}
\phantomsection
\addcontentsline{toc}{section}{\texorpdfstring{Exercises for \S\CurrentExerciseGroup}{Exercises for Section \CurrentExerciseGroup}}

\begin{enumerate}
\def\labelenumi{\arabic{enumi})}
\tightlist
\item
  Let H be a set, directed for the relation \(\leqslant\) , of integrable functions \(\geqslant 0\) such that \(\sup N_{1}(f) < +\infty\) . Let g be the upper envelope of H; in order that g be integrable and \(f \in H\)
\end{enumerate}

that, in \(L^{1}\) , the class \(\widetilde{g}\) of g be the limit of the filter of sections of the directed set of classes of the functions \(f \in H\) , it is necessary and sufficient that, for every \(\varepsilon > 0\) , there exist a function \(f_{1} \in H\) , a set B, directed for \(\leqslant\) , of lower semi-continuous integrable functions, and a mapping \(f \mapsto f^{*}\) of the set \(H_{1}\) of functions \(f \in H\) that are \(\leqslant f_{1}\) into the set B, such that \(f \leqslant f^{*}\) and \(\mathrm{N}_{1}(f^{*} - f) \leqslant \varepsilon\) for every function \(f \in H_{1}\) (make use of Th. 3 of No.~4).

\begin{enumerate}
\def\labelenumi{\arabic{enumi})}
\setcounter{enumi}{1}
\item
  Let \(\mu\) be Lebesgue measure on \(\mathbf{R}\) and let \(\Omega\) be the topological space obtained by equipping the space \(\mathcal{L}^1\) with the topology of pointwise convergence in \(\mathbf{R}\) . Show that the mapping \(f\mapsto \int fd\mu\) of \(\Omega\) into \(\mathbf{R}\) is not continuous at any point of \(\Omega\) .
\item
  Let I be an interval in \(\mathbf{R}\) , \(\mathbf{f}\) a mapping of \(X \times I\) into a Banach space \(F\) , such that: \(1^{\circ}\) for every \(\alpha \in I\) , the mapping \(t \mapsto \mathbf{f}(t, \alpha)\) of \(X\) into \(F\) is integrable; \(2^{\circ}\) for every \(t \in X\) , the mapping \(\alpha \mapsto \mathbf{f}(t, \alpha)\) has a derivative \(\mathbf{f}_{\alpha}'(t, \alpha)\) in \(I\) ; \(3^{\circ}\) there exists an integrable function \(g \geqslant 0\) such that \(|\mathbf{f}_{\alpha}'(t, \alpha)| \leqslant g(t)\) for all \(t \in X\) and all \(\alpha \in I\) . Under these conditions, show that the function \(\mathbf{u}(\alpha) = \int \mathbf{f}(t, \alpha) d\mu(t)\) is differentiable in \(I\) and that
\end{enumerate}

\[
\mathbf {u} ^ {\prime} (\alpha) = \int \mathbf {f} _ {\alpha} ^ {\prime} (t, \alpha) d \mu (t).
\]

\begin{enumerate}
\def\labelenumi{\arabic{enumi})}
\setcounter{enumi}{3}
\tightlist
\item
  Let \(\mu\) be Lebesgue measure on the interval \(X = [0,1]\) .
\end{enumerate}

\begin{enumerate}
\def\labelenumi{\alph{enumi})}
\item
  Define in X a nowhere dense perfect set A whose measure is an arbitrary number \(\alpha\) such that \(0 \leqslant \alpha < 1\) (use the method of construction of Cantor's triadic set).
\item
  Define in X a sequence \((\mathbf{A}_n)\) of pairwise disjoint nowhere dense integrable sets, such that \(\mu (\mathrm{A}_n) = 2^{-n}\) and such that every interval contiguous to \(\mathbf{B}_n = \bigcup_{k = 1}^{n}\mathbf{A}_k\) contains a subset of \(\mathrm{A}_{n + 1}\) of measure \(>0\) . If \(\mathbf{A} = \bigcup_{n}\mathbf{A}_n\) , show that \(\mathbf{A}\) is a meager set of measure 1, and \(\mathbb{C}\mathbf{A}\) a negligible non-meager set.
\item
  Let \(\mathrm{H} = \bigcup_{n=0}^{\infty} \mathrm{A}_{2n+1}\) ; show that, for every open interval \(\mathrm{I} \subset \mathrm{X}\) , the intersections of I with H and with CH have measure \textgreater{} 0. Let f be a function equal almost everywhere to the characteristic function \(\varphi_{H}\) ; show that there does not exist any sequence \((f_{n})\) of continuous functions on X, convergent at every point of X, whose limit is equal to f (note that f is necessarily discontinuous at every point of X, and make use of Exer. 22 of GT, IX, §5).
\end{enumerate}

\begin{enumerate}
\def\labelenumi{\arabic{enumi})}
\setcounter{enumi}{4}
\tightlist
\item
  Let \(\mu\) be a positive measure on a locally compact space \(X\) . For every numerical function \(f\) (finite or not), of whatever sign, defined on \(X\) , we denote by \(\mu^{*}(f)\) (and call upper integral of \(f\) ) the infimum of the numbers \(\mu(h)\) for the functions \(h \geqslant f\) that are integrable and lower semi-continuous, when such functions exist, and by \(+\infty\) in the contrary case; this definition coincides with that of §1, No.~3 when \(f \geqslant 0\) . We denote by \(\mu_{*}(f)\) , and call lower integral of \(f\) , the number \(-\mu^{*}(-f)\) .
\end{enumerate}

\begin{enumerate}
\def\labelenumi{\alph{enumi})}
\item
  Show that if \(f_{1}\) and \(f_{2}\) are two numerical functions such that \(f_{1}(x) \leqslant f_{2}(x)\) almost everywhere, then \(\mu^{*}(f_{1}) \leqslant \mu^{*}(f_{2})\) .
\item
  Let \(f_{1}\) and \(f_{2}\) be two numerical functions such that \(\mu^{*}(f_{1}) + \mu^{*}(f_{2})\) is defined and \(< +\infty\) ; show that \(f_{1}(x) + f_{2}(x)\) is defined almost everywhere and \(\mu^{*}(f_{1} + f_{2}) \leqslant \mu^{*}(f_{1}) + \mu^{*}(f_{2})\) (reduce to the case that \(f_{1}(x) < +\infty\) and \(f_{2}(x) < +\infty\) at every point, with the help of \(a\) )).
\item
  If \((f_n)\) is an increasing sequence of numerical functions such that \(\mu^{*}(f_{n}) > -\infty\) from some index onward, show that \(\mu^{*}(\sup_{n} f_{n}) = \sup_{n} \mu^{*}(f_{n})\) .
\end{enumerate}

¶ 6) a) Let \(f\) be a numerical function such that \(\mu^{*}(f)\) (Exer. 5) is finite. Show that there exists an integrable function \(f_{1} \geqslant f\) such that \(\mu(f_{1}) = \mu^{*}(f)\) ; if \(f_{2}\) is a second integrable function such that \(f_{2} \geqslant f\) and \(\mu(f_{2}) = \mu^{*}(f)\) , then \(f_{1}\) and \(f_{2}\) are equivalent.

\begin{enumerate}
\def\labelenumi{\alph{enumi})}
\setcounter{enumi}{1}
\item
  For a numerical function \(f\) to be integrable, it is necessary and sufficient that \(\mu^{*}(f)\) and \(\mu_{*}(f)\) be finite and equal.
\item
  Let \(f\) be a numerical function such that \(\mu^{*}(f)\) and \(\mu_{*}(f)\) are both finite; let \(g\) and \(h\) be two integrable functions such that \(g \leqslant f \leqslant h\) and \(\mu(g) = \mu_{*}(f)\) , \(\mu(h) = \mu^{*}(f)\) . Show that
\end{enumerate}

\[
\mu_ {*} (f - g) = \mu_ {*} (h - f) = 0 \quad \text { and } \quad \mu^ {*} (f - g) = \mu^ {*} (h - f) = \mu^ {*} (f) - \mu_ {*} (f).
\]

\begin{enumerate}
\def\labelenumi{\alph{enumi})}
\setcounter{enumi}{3}
\tightlist
\item
  Let \(f_{1}, f_{2}\) be two numerical functions such that the numbers \(\mu^{*}(f_{1}), \mu^{*}(f_{2})\) , \(\mu_{*}(f_{1})\) and \(\mu_{*}(f_{2})\) are all finite; show that
\end{enumerate}

\[
\mu_ {*} (f _ {1} + f _ {2}) \leqslant \mu_ {*} (f _ {1}) + \mu^ {*} (f _ {2}) \leqslant \mu^ {*} (f _ {1} + f _ {2})
\]

(if \(g_2\) is an integrable function such that \(f_2 \leqslant g_2\) and \(\mu^*(f_2) = \mu(g_2)\) , note that, for every integrable function \(h\) such that \(h \leqslant f_1 + f_2\) , one has \(h - g_2 \leqslant f_1\) ). From this, deduce that if \(f_1\) is integrable then

\[
\mu^ {*} (f _ {1} + f _ {2}) = \mu (f _ {1}) + \mu^ {*} (f _ {2}), \quad \mu_ {*} (f _ {1} + f _ {2}) = \mu (f _ {1}) + \mu_ {*} (f _ {2}),
\]

and \(\mu^{*}(f_{1} + f_{2}) = \mu^{*}\big(\sup (f_{1},f_{2})\big) + \mu^{*}\big(\inf (f_{1},f_{2})\big)\) (for the latter relation, make use of Exer. 1 of §1).

\begin{enumerate}
\def\labelenumi{\alph{enumi})}
\setcounter{enumi}{4}
\tightlist
\item
  Let \(f\) be an integrable function. In order that a function \(g\) , such that \(\mu^{*}(g)\) is finite, be integrable, it is necessary and sufficient that \(\mu(f) = \mu^{*}(g) + \mu^{*}(f - g)\) (if \(g_{1}\) is an integrable function such that \(g \leqslant g_{1}\) and \(\mu^{*}(g) = \mu(g_{1})\) , note that \(f - g_{1} \leqslant f - g\) ).
\end{enumerate}

¶ 7) Let \(\mu\) be a positive measure on X. For every subset \(A \subset X\) , the lower integral (Exer. 5) of the characteristic function \(\varphi_A\) is called the inner measure of A and is denoted \(\mu_*(A)\) .

\begin{enumerate}
\def\labelenumi{\alph{enumi})}
\item
  Show that \(\mu_{*}(\mathrm{A})\) is the supremum of the measures of the compact sets contained in A (argue as in Th. 4).
\item
  For every subset A of X of finite outer measure, show that there exist two integrable sets \(A_{1}, A_{2}\) such that \(A_{1} \subset A \subset A_{2}\) and \(\mu_{*}(A) = \mu(A_{1})\) , \(\mu^{*}(A) = \mu(A_{2})\) . For A to be integrable, it is necessary and sufficient that \(\mu^{*}(A)\) and \(\mu_{*}(A)\) be finite and equal. With the same notations, show that
\end{enumerate}

\[
\mu_ {*} (\mathrm{A} \cap \mathbb {C} \mathrm{A} _ {1}) = \mu_ {*} (\mathrm{A} _ {2} \cap \mathbb {C} \mathrm{A}) = 0
\]

and

\[
\mu^ {*} (\mathrm{A} \cap \mathbb {C} \mathrm{A} _ {1}) = \mu^ {*} (\mathrm{A} _ {2} \cap \mathbb {C} \mathrm{A}) = \mu^ {*} (\mathrm{A}) - \mu_ {*} (\mathrm{A}).
\]

\begin{enumerate}
\def\labelenumi{\alph{enumi})}
\setcounter{enumi}{2}
\item
  Let A be an integrable set; show that for every set \(\mathbf{B} \subset \mathbf{A}\) , one has \(\mu(\mathbf{A}) = \mu^{*}(\mathbf{B}) + \mu_{*}(\mathbf{A} \cap \mathbb{C}\mathbf{B})\) .
\item
  Let A and B be two sets of finite outer measure that are disjoint. Show that if \(\mathbf{C} = \mathbf{A} \cup \mathbf{B}\) , then
\end{enumerate}

\[
\mu_ {*} (\mathrm{A}) + \mu_ {*} (\mathrm{B}) \leqslant \mu_ {*} (\mathrm{C}) \leqslant \mu_ {*} (\mathrm{A}) + \mu^ {*} (\mathrm{B}) \leqslant \mu^ {*} (\mathrm{C}) \leqslant \mu^ {*} (\mathrm{A}) + \mu^ {*} (\mathrm{B})
\]

and that

\[
\mu_ {*} (\mathrm{C}) - \mu_ {*} (\mathrm{A}) - \mu_ {*} (\mathrm{B}) \leqslant \mu^ {*} (\mathrm{A}) + \mu^ {*} (\mathrm{B}) - \mu^ {*} (\mathrm{C})
\]

(for the latter inequality, reduce to the case that \(\mu_{*}(\mathrm{A}) = \mu_{*}(\mathrm{B}) = 0\) with the help of \(b\) ); if \(\mathrm{A}_2\) and \(\mathrm{B}_2\) are integrable sets such that \(\mathrm{A} \subset \mathrm{A}_2\) , \(\mathrm{B} \subset \mathrm{B}_2\) , \(\mu^{*}(\mathrm{A}) = \mu(\mathrm{A}_2)\) , \(\mu^{*}(\mathrm{B}) = \mu(\mathrm{B}_2)\) , show that \(\mu_{*}(\mathrm{C}) \leqslant \mu(\mathrm{A}_2 \cap \mathrm{B}_2)\) .

¶ 8) If the torus T = R/Z is canonically identified with the interval {[}0,1 {[} of R, the Lebesgue measure μ on this interval is a measure on T; for every set A ⊂ T and every z ∈ T, μ*(A + z) = μ*(A).

\begin{enumerate}
\def\labelenumi{\alph{enumi})}
\item
  Show that there exists a subgroup \(\mathrm{H}_0\) of \(\mathbf{T}\) such that the sets \(\mathrm{H}_n = r_n + \mathrm{H}_0\) , where \(r_n\) runs over the set of rational numbers contained in {[}0,1{[}, form a partition of \(\mathbf{T}\) (regarding \(\mathbf{R}\) as a vector space over \(\mathbf{Q}\) , note that in \(\mathbf{R}\) the subspace \(\mathbf{Q}\) admits a supplement).
\item
  Let H be a set that is the union of any finite number of sets \(H_{n}\) ; show that H is not integrable and that \(\mu_{*}(H)=0\) (note that every subgroup of the additive group Q/Z generated by a finite number of elements has infinite index in Q/Z; from this, deduce that there exists a partition of T into a countable infinity of sets \(P_{n}\) , each of which contains a set of the form \(z+H\) ).
\item
  Deduce from \(b\) ) an example of a decreasing sequence \((\mathrm{A}_n)\) of subsets of \(\mathbf{T}\) , whose intersection is empty, and is such that \(\mu^{*}(\mathrm{A}_{n}) = 1\) for all \(n\) .
\item
  Let A be an integrable set such that \(H_{0} \subset A\) and \(\mu(A) = \mu^{*}(H_{0}) > 0\) ; show that, for every integrable set \(B \subset A\) , the sets \(B_{1} = B \cap H_{0}\) and \(B_{2} = B \cap C H_{0}\) form a partition of B such that \(\mu^{*}(B_{1}) = \mu^{*}(B_{2}) = \mu(B)\) and \(\mu_{*}(B_{1}) = \mu_{*}(B_{2}) = 0\) (make use of Exer. 7 b)).
\end{enumerate}

¶ 9) a) Let \(\mu\) be a positive measure on a locally compact space \(X\) . In the Banach space \(\widetilde{\mathcal{F}}^1\) of equivalence classes of functions in \(\mathcal{F}^1\) , let \(G\) be a closed linear subspace containing the subspace \(L^{1}\) ; let G be the closed linear subspace of \(F^{1}\) formed by the functions whose class belongs to G. Show that one can extend to G the linear form \(\mu(f)\) in such a way that the inequality \(|\mu(f)| \leqslant N_{1}(f)\) is again verified (make use of the Hahn--Banach theorem). Show that for the integral so extended, Th. 3 of §3 is again valid.

\begin{enumerate}
\def\labelenumi{\alph{enumi})}
\setcounter{enumi}{1}
\tightlist
\item
  Assume that G is the direct sum of \(\mathbf{L}^1\) and a subspace H of finite dimension. Show that, in \(\mathcal{G}\) , Lebesgue's theorem (Th. 2) is again valid. (Let \((f_n)\) be a sequence of functions in \(\mathcal{G}\) , tending almost everywhere to \(f\) and such that \(|f_n| \leqslant g\) , where \(g \geqslant 0\) and \(\mathrm{N}_1(g) < +\infty\) . Let \((\widetilde{u}_k)_{1 \leqslant k \leqslant m}\) be a basis of H, and let \(f_n = \sum_{k=1}^{m} \alpha_{nk} u_k + h_n\) , where \(h_n \in \mathcal{L}^1\) . Using the fact that \(G\) is the topological direct sum of \(L^1\) and \(H\) , show that the \(\alpha_{nk}\) are uniformly bounded; on passing, if necessary, to a subsequence of \((f_n)\) , reduce to the case that each of the sequences \((\alpha_{nk})_{n \geqslant 1}\) has a limit; from this, deduce that \(h_n(x)\) then tends to a limit almost everywhere, and apply Lebesgue's theorem to the sequence \((h_n)\) ; note finally that two subsequences of \((\widetilde{f}_n)\) cannot tend to different limits in \(G\) .
\end{enumerate}

¶ 10) a) Let E be a Riesz space, μ a positive linear form on E such that the relation μ(\textbar x\textbar) = 0 implies x = 0; μ(\textbar x\textbar) is then a norm on E, and let us assume that E, equipped with this norm, is complete. It then follows that E is fully lattice-ordered (Ch. II, §2, Exer. 8 e)). Consequently (Ch. II, §1, Exers. 12 and 13) there exists a locally compact space X, the sum of a family (Kα) of compact Stone spaces, such that E is isomorphic to a space of continuous numerical functions (finite or not) on X that contains the space K(X). Identifying E with this space, the restriction of μ to K(X) is then a positive measure on X. Show that E is canonically isomorphic to the space L¹(μ); more precisely, for every μ-integrable function g defined on X, there exists one and only one function f ∈ E that is equivalent to g for the measure μ (note that every element ≥ 0 of E is the supremum of an increasing sequence of elements of K(X), and that K(X) is dense in L¹(μ)).

\begin{enumerate}
\def\labelenumi{\alph{enumi})}
\setcounter{enumi}{1}
\tightlist
\item
  Deduce from a) that, for every compact space K, there exist a locally compact space S, the topological sum of a family of compact Stone spaces, and a positive measure \(\nu\) on S, with support equal to S, such that the fully lattice-ordered space \(\mathcal{M}(\mathrm{K})\) of measures on K, equipped with the norm \(\|\mu\|\) , is isomorphic to \(\mathcal{L}^{1}(\nu)\) (consider the linear form \(\mu\mapsto\mu(\mathrm{K})\) on \(\mathcal{M}(\mathrm{K})\) ).
\end{enumerate}

\begin{enumerate}
\def\labelenumi{\arabic{enumi})}
\setcounter{enumi}{10}
\tightlist
\item
  \begin{enumerate}
  \def\labelenumii{\alph{enumii})}
  \tightlist
  \item
    Let \(\Gamma\) be any set of subsets of a set \(A\) . Let \(\Psi\) the set of subsets of \(A\) of the form
  \end{enumerate}
\end{enumerate}

\[
\mathrm{X} _ {1} \cap \mathrm{X} _ {2} \cap \dots \cap \mathrm{X} _ {m} \cap \mathbb {C} \mathrm{X} _ {m + 1} \cap \dots \cap \mathbb {C} \mathrm{X} _ {m + p},
\]

where the \(X_{i}\) are sets in \(\Gamma\) , m is any integer \(\geqslant 1\) , and p is any integer \(\geqslant 0\) . Show that the smallest clan \(\Phi\) containing \(\Gamma\) is the set of finite unions of sets in \(\Psi\) .

\begin{enumerate}
\def\labelenumi{\alph{enumi})}
\setcounter{enumi}{1}
\item
  Let \(\Delta\) be the set of finite intersections of sets in \(\Gamma\) . Show that, for every vector space F, the set of linear combinations (with coefficients in F) of characteristic functions of sets of the clan generated by \(\Gamma\) is identical to the set of linear combinations of characteristic functions of sets in \(\Delta\) .
\item
  Let X be a topological space, \(\Gamma\) the set of compact subsets of X. Show that the clan generated by \(\Gamma\) is identical to the set of finite unions of subsets of X of the form \(A \cap CB\) , where A and B are compact sets.
\end{enumerate}

\begin{enumerate}
\def\labelenumi{\arabic{enumi})}
\setcounter{enumi}{11}
\tightlist
\item
  Let \(X\) be a Hausdorff topological space. For every pair \((K, U)\) formed by a compact set \(K\) and an open set \(U\) in \(X\) , let \(I(K, U)\) be the set of subsets \(M \subset X\) such that \(K \subset M \subset U\) , and let \(\mathcal{T}\) be the topology on \(\mathfrak{P}(X)\) generated by the set of subsets \(I(K, U)\) of \(\mathfrak{P}(X)\) .
\end{enumerate}

\begin{enumerate}
\def\labelenumi{\alph{enumi})}
\item
  Show that each of the sets I(K, U) is both open and closed in \(\mathfrak{P}(X)\) ; from this, deduce that \(\mathfrak{P}(X)\) , equipped with the topology \(\mathcal{T}\) , is a completely regular and totally disconnected space.
\item
  For the mapping \(M \mapsto CM\) of \(\mathfrak{P}(X)\) onto itself to be continuous (for the topology T), it is necessary and sufficient that X be compact.
\item
  Take X to be the interval {[}0,1{]} of R. Show that the mappings \((M,N)\mapsto M\cup N\) and \((M,N)\mapsto M\cap N\) of \(\mathfrak{P}(X)\times\mathfrak{P}(X)\) into \(\mathfrak{P}(X)\) are not continuous for the topology T.
\item
  Assume X to be locally compact. Show that the topology induced by T on the set of compact subsets of X is finer than the topology deduced from a uniform structure on X by the procedure of Exer. 5 of GT, II, §1; these two topologies cannot be identical unless X is a discrete space.
\end{enumerate}

¶ 13) a) Let X be a locally compact space, and let \(\Gamma\) be a base for the topology of X consisting of relatively compact sets. Let \(\mathcal{U}\) be a uniform structure compatible with the topology of X, and let \(\mathfrak{S}\) be a fundamental system of entourages for this structure. Let \(M \mapsto \lambda(M)\) be a finite numerical function \(\geqslant 0\) defined on \(\Gamma\) . For every compact set \(K \subset X\) and every entourage \(V \in \mathfrak{S}\) , let \(\alpha_V(K)\) be the infimum of the numbers \(\sum_{i} \lambda(U_i)\)

for all finite coverings \((\mathrm{U}_{i})\) of K formed by sets in \(\Gamma\) that are small of order V; suppose that as V runs over S, the supremum \(\alpha(\mathrm{K})\) of the numbers \(\alpha_{\mathrm{V}}(\mathrm{K})\) is finite. Show that there exists a measure \(\mu\) on X (and only one) such that \(\mu(\mathrm{K}) = \alpha(\mathrm{K})\) for every compact set K (make use of Th. 5).

\begin{enumerate}
\def\labelenumi{\alph{enumi})}
\setcounter{enumi}{1}
\tightlist
\item
  Let \(\Psi\) be the set of Borel subsets of \(X\) , \(\alpha\) a mapping of \(\Psi\) into \([0, +\infty]\) satisfying the following conditions:
\end{enumerate}

\begin{enumerate}
\def\labelenumi{(\roman{enumi})}
\item
  If \(\mathrm{B}_1, \mathrm{B}_2\) are two disjoint Borel subsets of \(X\) , then \(\alpha(\mathrm{B}_1 \cup \mathrm{B}_2) = \alpha(\mathrm{B}_1) + \alpha(\mathrm{B}_2)\) .\\
\item
  If \(B\) is a compact subset of \(X\) then \(\alpha(B) < +\infty\) .
\item
  If \(\mathbf{B}\) is a Borel subset of \(\mathbf{X}\) then \(\alpha(\mathbf{B}) = \inf \alpha(\mathbf{U})\) , where \(\mathbf{U}\) runs over the set of open subsets of \(\mathbf{X}\) containing \(\mathbf{B}\) .
\end{enumerate}

Then, there exists one and only one positive measure \(\mu\) on X such that \(\alpha(B) = \mu^{*}(B)\) for every Borel subset B of X that can be covered by a sequence of compact sets. c) For every \(B \in \Phi\) , set \(\alpha(B) = 0\) if B can be covered by a sequence of compact sets, and \(\alpha(B) = +\infty\) in the contrary case. Then, the conditions (i), (ii), (iii) of \(b\) are satisfied. One has \(\mu = 0\) , therefore \(\alpha(B) \neq \mu^{*}(B)\) if B cannot be covered by a sequence of compact sets.

\begin{enumerate}
\def\labelenumi{\arabic{enumi})}
\setcounter{enumi}{13}
\item
  Let \((\mathbf{A}_n)\) be a sequence of integrable sets such that \(\sum_{n} \mu(\mathbf{A}_n) < +\infty\) . For every integer \(k\) , let \(\mathbf{G}_k\) be the set of \(x \in X\) such that \(x \in \mathbf{A}_n\) for at least \(k\) values of \(n\) ; show that \(\mathbf{G}_k\) is integrable and that \(k \cdot \mu(\mathbf{G}_k) \leqslant \sum_{n} \mu(\mathbf{A}_n)\) .
\item
  Let \(\mu\) be a bounded positive measure on a locally compact space \(X\) , and let \((A_n)\) be a sequence of integrable sets in \(X\) such that \(\inf_{n} \mu(A_n) = m > 0\) . Show that the set \(B\) of points of \(X\) that belong to infinitely many of the sets \(A_n\) is integrable, and that \(\mu(B) \geqslant m\) .
\end{enumerate}

¶ 16) Let X be a completely regular space, and let \(\mathcal{C}(X)\) (resp. \(\mathcal{C}^b(X)\) ) be the Riesz space of continuous (resp. continuous and bounded) numerical functions on X.

\begin{enumerate}
\def\labelenumi{\alph{enumi})}
\item
  Show that if a linear form \(\lambda\) on \(\mathcal{C}(\mathrm{X})\) is continuous for the topology of compact convergence, it is relatively bounded.
\item
  Let \(\lambda\) be a positive linear form on \(\mathcal{C}(\mathrm{X})\) . Show that if \(\lambda\) is zero on \(\mathcal{C}^b (\mathrm{X})\) , then it is zero on \(\mathcal{C}(\mathrm{X})\) (let \(\varphi\) be the canonical mapping of \(\mathcal{C}(\mathrm{X})\) onto the quotient Riesz space \(\mathcal{C}(\mathrm{X}) / \mathcal{C}^b (\mathrm{X})\) (Ch. II, §1, Exer. 4); show that if \(h\) is a continuous function \(\geqslant 0\) and is not bounded on \(\mathrm{X}\) , then \(n\varphi (h)\leqslant \varphi (h^{2})\) for every integer \(n > 0\) ).
\item
  Let \(\beta X\) be the Stone--Čech compactification of \(X\) , the compact space obtained by equipping \(X\) with the coarsest uniform structure for which the functions in \(\mathcal{C}^b(X)\) are uniformly continuous and by completing the uniform space so obtained; every function \(f \in \mathcal{C}(X)\) may then be extended by continuity to a continuous function \(\widetilde{f}\) on \(\beta X\) , with possibly infinite values (consider \(f/(1 + |f|)\) ). If \(\lambda\) is a positive linear form on \(\mathcal{C}(X)\) , then the restriction of \(\lambda\) to \(\mathcal{C}^b (\mathrm{X})\) is of the form \(f\mapsto \mu (\widetilde{f})\) , where \(\mu\) is a positive measure on \(\beta X\) ; show that for every function \(f\in \mathcal{C}(\mathrm{X})\) , \(\widetilde{f}\) is integrable for \(\mu\) and \(\lambda (f) = \mu (\widetilde{f})\) (make use of \(b\) ), noting that every function \(\geqslant 0\) belonging to \(\mathcal{C}(\mathrm{X})\) is the upper envelope of a sequence of functions in \(\mathcal{C}^b (\mathrm{X})\) ).
\item
  Show that every point \(x_{0}\) of the support of \(\mu\) that does not belong to X has the following property: for every decreasing sequence \((V_{n})\) of neighborhoods of \(x_{0}\) , the intersection of the \(V_{n}\) contains at least one point of X. Converse.
\item
  Deduce from \(d\) ) that if \(X\) is locally compact and countable at infinity, then the support of \(\mu\) is contained in \(X\) (compare with \(h\) )).
\item
  If X is discrete, show that the support of \(\mu\) is finite (in the contrary case, form a function \(f \geqslant 0\) defined on X such that \(\widetilde{f}\) is not \(\mu\) -integrable).
\item
  Show that if a linear form \(\lambda\) on \(\mathcal{C}(\mathbf{X})\) is positive and continuous for the topology of compact convergence, then the support of \(\mu\) is contained in X (cf.~Ch. III, §2, No.~3, Prop. 11).
\item
  Let \(X_{0}\) be the compact space obtained by adjoining a point at infinity \(\omega\) to the locally compact space X defined in GT, I, §9, Exer. 11, let Y be the subspace of \(\overline{R}\) formed by the integers \(\geqslant 0\) and \(+\infty\) , and let Z be the locally compact space complementary to the point \((\omega, +\infty)\) in the product space \(X_{0} \times Y\) . Show that every continuous numerical function on Z is bounded and that \(f(z)\) tends to a limit as z tends to the point at infinity \((\omega, +\infty)\) of Z. From this, deduce that if \(\mu\) is the measure on Z defined by the mass \(1/2^{n}\) placed at the point \((\omega, n)\) for every \(n \geqslant 0\) , then every continuous function on Z is \(\mu\) -integrable, but the support of \(\mu\) is not compact.
\end{enumerate}

¶ 17) Let X be the locally compact space defined in GT, I, §9, Exer. 11.

\begin{enumerate}
\def\labelenumi{\alph{enumi})}
\tightlist
\item
  Let H be a compact subset of the locally convex space \(\mathcal{C}(\mathrm{X};\mathbf{C})\) equipped with the topology of compact convergence; show that the functions \(f \in H\) are uniformly bounded on X and that there exists \(c \in X\) such that, for every \(x \geqslant c\) , all of the functions of H are constant. (Argue by contradiction; if the functions of H were not uniformly bounded on X, there would be an increasing sequence \((x_{n})\) of points of X such that \(\sup(|\mathrm{H}(x_{n})|) \geqslant n\) , where \(|\mathrm{H}(x_{n})|\) denotes the set of \(|f(x_{n})|\) for \(f \in H\) ; observe that the sequence \((x_{n})\) is convergent in X. Similarly, for every \(x \in X\) , let \(\delta(x)\) be the supremum of the oscillations of all the functions \(f \in H\) in the interval \([x, \rightarrow[; \delta(x)\) is finite and decreasing, therefore there exists \(d \in X\) such that \(\delta(x)\) is constant for \(x \geqslant d\) . To see that this constant \(\beta\) is zero, argue by contradiction as before, on forming two increasing sequences \((s_{n}), (t_{n})\) of points of X such that \(s_{n} \leqslant t_{n} \leqslant s_{n+1} \leqslant t_{n+1}\) and a sequence \((f_{n})\) of functions in H such that \(|f_{n}(s_{n}) - f_{n}(t_{n})| \geqslant \beta/2\) .
\end{enumerate}

In particular, \(\mathcal{C}(\mathrm{X};\mathbf{C})\) is the direct sum of \(\mathcal{K}(\mathrm{X};\mathbf{C})\) and \(C\cdot1\) .

\begin{enumerate}
\def\labelenumi{\alph{enumi})}
\setcounter{enumi}{1}
\item
  Show that on X, every measure has compact support (for every \(x \in X\) let \(f_x\) be the characteristic function of the interval \(\leftarrow, x\) {]}, which is continuous and has compact support; consider on X the increasing function \(|\mu|(f_x)\) ).
\item
  Deduce from a) and b) that in \(\mathcal{M}(\mathrm{X};\mathbf{C}) = \mathcal{M}^{1}(\mathrm{X};\mathbf{C}) = \mathcal{C}'(\mathrm{X};\mathbf{C})\) , the bounded sets for the topology \(\mathcal{T}\) of uniform convergence on the compact subsets of \(\mathcal{C}(\mathrm{X};\mathbf{C})\) are the bounded sets for the ultrastrong topology (Ch. III, §1, Exer. 15). Show that \(\mathcal{C}'(\mathrm{X};\mathbf{C})\) is not quasi-complete for the topology \(\mathcal{T}\) (consider the measures \(\varepsilon_{x}\) for \(x\in \mathrm{X}\) ); from this, deduce that for the topology of compact convergence, \(\mathcal{C}(\mathrm{X};\mathbf{C})\) is neither barreled nor bornological (cf.~TVS, III, §4, No.~2, Cor. 4 of Th. 1 and Exer. 18 below).
\item
  Show that on \(\mathcal{K}(\mathrm{X};\mathbf{C})\) the topology \(\mathcal{T}_0\) (the direct limit of the topologies of the Banach spaces \(\mathcal{K}(\mathrm{X},\mathrm{K};\mathbf{C})\) ) is identical to the topology of uniform convergence, and that, equipped with this topology, \(\mathcal{K}(\mathrm{X};\mathbf{C})\) is complete. (Let V be a neighborhood of 0 for \(\mathcal{T}_0\) ; for every \(x\in X\) , let \(r_x\) be the largest number \(>0\) such that V contains all of the continuous functions \(f\) with support contained in \([ \leftarrow ,x]\) and such that \(\| f\| \leqslant r_x\) . Show that the infimum of \(r_x\) in X is \(>0\) , by proving that there cannot exist an increasing sequence \((x_n)\) of points of X such that \(\lim_{n\to \infty}r_{x_n} = 0\) .)
\end{enumerate}

\begin{enumerate}
\def\labelenumi{\arabic{enumi})}
\setcounter{enumi}{17}
\tightlist
\item
  Let \(X\) be a paracompact locally compact space.
\end{enumerate}

\begin{enumerate}
\def\labelenumi{\alph{enumi})}
\item
  Show that the space \(\mathcal{C}(\mathrm{X};\mathbf{C})\) is isomorphic to a product of Fréchet spaces, hence is barreled.
\item
  For a subset H of \(\mathcal{C}^{\prime}(X; \mathbf{C})\) to be bounded for the topology of compact convergence, it is necessary and sufficient that there exist a compact subset K of X such that \(\operatorname{Supp}(\mu) \subset K\) for all \(\mu \in H\) and such that H is bounded for the ultrastrong topology (compare with Exer. 17 c)). The topology induced on H by the topology of compact convergence is then identical to the vague topology.
\item
  Show that on the set \(\mathcal{M}_{+}(\mathrm{X})\cap \mathcal{C}'(\mathrm{X};\mathbf{C})\) the topology induced by the topology of compact convergence on \(\mathcal{C}'(\mathrm{X};\mathbf{C})\) is identical to the vague topology. Is the same true when \(X\) is not paracompact (Exer. 17)?
\end{enumerate}

\begin{enumerate}
\def\labelenumi{\arabic{enumi})}
\setcounter{enumi}{18}
\tightlist
\item
  Equip \(\mathcal{C}(\mathrm{X};\mathbf{C})\) with the topology of compact convergence, and \(\mathcal{C}'(\mathrm{X};\mathbf{C})\) with the topology of uniform convergence on the compact subsets of \(\mathcal{C}(\mathrm{X};\mathbf{C})\) . Let \(\mathfrak{S}\) be the set of compact subsets of \(\mathcal{C}(\mathrm{X};\mathbf{C})\) ; show that the bilinear mapping \((g,\mu)\mapsto g\cdot \mu\) of \(\mathcal{C}(\mathrm{X};\mathbf{C})\times \mathcal{C}'(\mathrm{X};\mathbf{C})\) into \(\mathcal{C}'(\mathrm{X};\mathbf{C})\) is \(\mathfrak{S}\) -hypocontinuous. Let \(\mathcal{T}\) be the set of bounded subsets of \(\mathcal{C}'(\mathrm{X};\mathbf{C})\) ; show that if X is paracompact, the preceding bilinear mapping is also \(\mathcal{T}\) -hypocontinuous; is the latter property again true when X is no longer assumed to be paracompact (Exer. 17)?
\end{enumerate}

¶ 20) Let X,Y be two paracompact locally compact spaces. Equip the spaces \(\mathcal{C}'(X; \mathbf{C})\) and \(\mathcal{C}'(Y; \mathbf{C})\) with the topology of compact convergence.

\begin{enumerate}
\def\labelenumi{\alph{enumi})}
\item
  Show that when \(\mathcal{C}^{\prime}(X\times Y;\mathbf{C})\) is equipped with the topology of compact convergence, the mapping \((\mu,\nu)\mapsto\mu\otimes\nu\) is \((\mathfrak{S},\mathfrak{T})\) -hypocontinuous, where S denotes the set of bounded subsets of \(\mathcal{C}^{\prime}(X;\mathbf{C})\) and T the set of bounded subsets of \(\mathcal{C}^{\prime}(Y;\mathbf{C})\) (make use of Exer. 18 b)).
\item
  Show that when X and Y are countable at infinity and \(\mathcal{C}^{\prime}(X\times Y;C)\) is equipped with the topology of compact convergence, the mapping \((\mu,\nu)\mapsto\mu\otimes\nu\) is continuous. (Inspired by the proof of Exer. 3 of Ch. III, §4 and introducing in X (resp. Y) a continuous partition of unity subordinate to a locally finite open covering, argue as in TVS, II, §4, Exer. 9 a).
\item
  If X is discrete, the space \(\mathcal{C}^{\prime}(X;\mathbf{C})\) may be identified with the direct sum space \(\mathbf{C}^{(X)}\) , equipped with the finest locally convex topology. From this, deduce an example where X and Y are discrete, with X uncountable, and where the mapping \((\mu,\nu)\mapsto\mu\otimes\nu\) is not continuous when \(\mathcal{C}^{\prime}(X\times Y;\mathbf{C})\) is equipped with the vague topology (cf.~TVS, II, §4, Exer. 9 b)).
\end{enumerate}

\begin{enumerate}
\def\labelenumi{\arabic{enumi})}
\setcounter{enumi}{20}
\item
  Let \(\mathbf{X}\) be a locally compact space, \(f\) a numerical function \(\geqslant 0\) defined on \(\mathbf{X}\) . In order that the mapping \(\mu \mapsto \mu^{*}(f)\) of \(\mathcal{M}_{+}(\mathbf{X}) \cap \mathcal{C}'(\mathbf{X};\mathbf{C})\) into \(\overline{\mathbf{R}}\) be continuous for the topology of compact convergence, it is necessary and sufficient that \(f\) be continuous on \(\mathbf{X}\) .
\item
  Let X, Y be two locally compact spaces, \(\mu\) a bounded measure on X, \(\nu\) a bounded measure on Y, so that \(\mu \otimes \nu\) is a bounded measure on \(X \times Y\) . Show that for every function \(f \in \mathcal{C}^{0}(X \times Y; \mathbf{C})\) , the function \(x \mapsto \int f(x, y) d\nu(y)\) belongs to \(\mathcal{C}^{0}(X; \mathbf{C})\) and
\end{enumerate}

\[
\int f (x, y) d \mu (x) d \nu (y) = \int d \mu (x) \int f (x, y) d \nu (y).
\]

¶ 23) On the space \(\mathbf{R}^n\) , let \(\mu\) be Lebesgue measure, \(|\mathbf{x}|\) a norm such that the unit ball for this norm has measure equal to 1, and \((\mathbf{x}_k)_{k \geqslant 1}\) an infinite sequence of distinct points of a bounded integrable set B such that \(\mu(B) = 1\) . For every integer \(m\) , denote by \(d_m\) the infimum of the numbers \(|\mathbf{x}_i - \mathbf{x}_j|\) for \(1 \leqslant i < j \leqslant m\) . Show that

\(\liminf_{m \to \infty} md_m^n \leqslant \alpha_n^{-1}\) , where

\[
\alpha_ {n} = 1 + n \int_ {0} ^ {1} \frac {(1 - t) ^ {n}}{1 + t} d t.
\]

(Argue by contradiction, assuming that for some \(\varepsilon > 0\) there exists an \(m_0\) such that \(md_m^n > h_n\) for \(m \geqslant m_0\) , with \(h_n = \alpha_n^{-1} + \varepsilon\) . For \(1 \leqslant i < m\) , let \(B_i\) be the ball with center \(\mathbf{x}_i\) and radius \(\frac{1}{2} h_n m^{-1/n}\) , and for \(m \leqslant i \leqslant 2^n m\) let \(B_i\) be the ball with center \(\mathbf{x}_i\) and radius \(\frac{1}{2} h_n (2i^{-1/n} - m^{-1/n})\) . Show that the \(2^n m\) balls \(B_i\) are pairwise disjoint, and evaluate the measure of their union, making use of the Euler--Maclaurin formula.)
